\chapter{分析力学与几何力学：Lagrange--Hamilton 方程、Euler--Poincaré 与辛结构}
\label{ch:geometric-mechanics}

% ════════════════════════════════════════════════════════════════
%  控制部 C2 第 K 章（卷五《机器人最优控制》刚体动力学子部第二章）。
%  吸收 Robotics_Tutorial 50d/20(Lagrange-Hamilton) + 50d/40(SE(3)几何力学)
%  + 50d/70(辛结构与动量映射)，研究生密度，纯理论（范围 A，库/代码全删）。
%  与 J 章(ch:spatial-dynamics) 强配套：J=计算视角，K=分析力学/几何视角。
%  6D 全章 [omega;v] 约定，与本书 se(3) ξ=[rho;φ] 显式置换对照；
%  ad* 用 math 约定 ad*=ad^T，对接 Featherstone force cross 处标 ×*=−ad^T。
%  勘误 E1–E7 全部已应用（设计稿§2 + 复审裁决：6 真错 1 部分）。
% ════════════════════════════════════════════════════════════════

\begin{note}
\textbf{本章定位与记号约定.}\;
导论 \cref{ch:control} 已给出计算力矩控制、阻抗、WBC 的全景与 $M\ddot{\mathbf q}+\mathbf C\dot{\mathbf q}+\mathbf g=\boldsymbol\tau$ 的工程用法；姊妹章 \cref{ch:spatial-dynamics} 从\emph{计算视角}用 6D 空间向量与 $O(n)$ 递推（RNEA/ABA）把这个方程算了出来。本章从\emph{分析力学与几何视角}回答它\emph{从何而来}：把同一个方程升级为「变分原理／李群约化／辛流」三条殊途同归的几何叙事。三部分依次为分析力学（\crefrange{sec:gm-history}{sec:gm-gravity}：从 Hamilton 原理到 $\mathbf M,\mathbf C,\mathbf g$ 与 $\dot{\mathbf M}-2\mathbf C$ 反对称）、几何力学（\crefrange{sec:gm-why-geom}{sec:gm-centroidal}：Euler--Poincaré、浮基、Lie--Poisson）、辛结构（\crefrange{sec:gm-symplectic}{sec:gm-frontier}：辛流形、动量映射、Marsden--Weinstein 约化、辛积分器）。

\textbf{记号统一.}\;旋转 $\mathbf R\in\SO(3)$、位姿 $\mathbf T\in\SE(3)$，伴随／余伴随／指数对数沿用 \cref{ch:lie}。\textbf{6D 空间向量全章采角在前序 $[\boldsymbol\omega;\mathbf v]$}（几何力学／Featherstone／Lynch--Park 主流），与本书 $\se(3)$ 坐标 $\boldsymbol\xi=[\boldsymbol\rho;\boldsymbol\phi]$（线在前）相反，二者由置换 $\mathbf P=\bigl(\begin{smallmatrix}\mathbf 0&\mathbf I\\\mathbf I&\mathbf 0\end{smallmatrix}\bigr)$ 互换（\cref{sec:gm-bodyvel} 明示）。\textbf{余伴随采数学约定 $\operatorname{ad}^*=\operatorname{ad}^\top$}（与 \cref{ch:lie} 一致），仅在对接 \cref{ch:spatial-dynamics} 的 RNEA 力回溯（force cross）处标注 $\times^*=-\operatorname{ad}^*$。$6\times6$ 空间惯性记 $\boldsymbol{\mathcal I}$（花体，与单位阵 $\mathbf I$、$3\times3$ 转动惯量 $\mathbf I_c$ 区分）。辛标准矩阵 $\mathbf J=\bigl(\begin{smallmatrix}\mathbf 0&\mathbf I\\-\mathbf I&\mathbf 0\end{smallmatrix}\bigr)$ 仅在辛节局部使用，动量映射用正体 $\mathsf J$ 与之区分，二者皆不与 \cref{ch:lie} 的右雅可比 $\mathbf J_r$、运动学 Jacobian 混淆。广义动量三身份：关节共轭动量 $\mathbf p=\mathbf M\dot{\mathbf q}$（$n$ 维）、李代数对偶元 $\boldsymbol\mu=\partial\ell/\partial\boldsymbol\xi$（$6$ 维）、体角动量 $\boldsymbol\pi=\mathbf I_c\boldsymbol\omega$（$3$ 维）。
\end{note}

% ════════════════════════════════════════════════════════════════
% 第一部分 — 分析力学
% ════════════════════════════════════════════════════════════════
\section{动机与历史：一个方程统治控制、仿真与学习}
\label{sec:gm-history}

\paragraph{动机：标准方程是物理与优化之间的桥。}
现代机器人控制、仿真、学习共享同一条运动方程
\begin{equation}
  \mathbf M(\mathbf q)\ddot{\mathbf q}+\mathbf C(\mathbf q,\dot{\mathbf q})\dot{\mathbf q}+\mathbf g(\mathbf q)=\boldsymbol\tau,
  \label{eq:gm-standard}
\end{equation}
其中 $\mathbf q\in\mathcal Q$ 为广义坐标、$\mathbf M$ 为质量（惯性）矩阵、$\mathbf C\dot{\mathbf q}$ 为 Coriolis 与离心力、$\mathbf g$ 为重力、$\boldsymbol\tau$ 为广义力。计算力矩控制用它做反馈线性化 $\boldsymbol\tau=\mathbf M\ddot{\mathbf q}_d+\mathbf C\dot{\mathbf q}+\mathbf g$；MPC 把它当等式约束；仿真器解正动力学 $\ddot{\mathbf q}=\mathbf M^{-1}(\boldsymbol\tau-\mathbf C\dot{\mathbf q}-\mathbf g)$；系统辨识估计其惯性参数；基于模型的强化学习把它当环境模型。会推导 \cref{eq:gm-standard} 并理解其结构，等于把「机器人是二阶受控动力系统」落实成可求解、可分析的数学对象。

\paragraph{反面：只会 Newton--Euler 会怎样。}
Newton--Euler（即 \cref{ch:spatial-dynamics} 的 RNEA）能高效算出力矩，却不直接给出 \cref{eq:gm-standard} 的\emph{结构}。而控制器设计恰恰依赖结构：被动性与自适应控制的 Lyapunov 证明需要 $\dot{\mathbf M}-2\mathbf C$ 的反对称性（\cref{thm:gm-skew}），这一性质只在 Lagrange 框架下自然显现；Hamilton 力学给出能量守恒与辛结构，决定长时仿真该选哪种积分器；Noether 定理由对称性\emph{直接}读出守恒量（浮基动量守恒、旋转对称的角动量守恒）。换言之，Lagrange/Hamilton 框架提供的不是「另一种算法」，而是「方程为什么长这样」的解释力——本章正是要把这套解释力建立起来。

\paragraph{历史脉络。}
分析力学的主线可上溯至 d'Alembert（1743，虚功原理把动力学化为静力学）、Lagrange（1788，《分析力学》以广义坐标统一力学）、Hamilton（1834--1835，「动力学的一般方法」引入正则形式与最小作用量原理）、Jacobi（1837，正则变换理论将 Hamilton 体系推广为 Hamilton--Jacobi 方法）、Lie（1870 年代，连续对称群）、Noether（1918，对称性与守恒量的对应）、Arnold（1960 年代，辛几何的现代形式化）\footnote{常见史料把 Hamilton 原理列为 1833、Hamilton--Jacobi 列为 1835；据 Nakane--Fraser 对早期 Hamilton--Jacobi 动力学的考证，Hamilton 的奠基文 \emph{On a General Method in Dynamics} 实为 1834（次年续篇引入正则形式），Jacobi 的推广在 1836--1837，本书采订正后的年份。}。机器人学侧的两座里程碑是 Walker--Orin（1982，用 RNEA 构造 $\mathbf M$ 的 CRBA 算法，见 \cref{ch:spatial-dynamics}）与 Slotine--Li（1987，利用 $\dot{\mathbf M}-2\mathbf C$ 反对称设计自适应控制器\cite{slotine1987adaptive}）。本章把这条主线中「分析力学—几何力学—辛」的理论内核完整重建，并指明其每一处的机器人接口。

% ════════════════════════════════════════════════════════════════
\section{广义坐标与约束}
\label{sec:gm-coords}

\paragraph{位形空间与广义坐标。}
$n$ 自由度系统的\emph{位形空间} $\mathcal Q$ 是所有可能构型之集，$\mathbf q=(q_1,\dots,q_n)\in\mathcal Q$，每个 $q_i$ 为一个\emph{广义坐标}——可以是角度、弧长或任何方便的参数，不必是笛卡尔坐标。用广义坐标的根本好处是\emph{自动满足约束}：一个 $n$ 连杆串联臂若以各连杆的笛卡尔位姿描述需 $6n$ 个变量加关节约束方程，而以 $n$ 个关节角描述则只用 $n$ 个变量——关节约束已被坐标选择吸收。平面 $n$R 臂位形空间为环面 $\mathbb T^n$；浮基系统（四足、人形）则天然是 $\mathcal Q=\SE(3)\times\mathbb T^n$，这一点将在 \cref{sec:gm-why-geom} 成为几何力学的出发点。

\begin{definition}[完整约束与非完整约束]\label{def:gm-constraint}
可写成位置级等式 $\boldsymbol\phi(\mathbf q,t)=\mathbf 0$ 的约束称为\emph{完整}（holonomic）约束，它把位形空间从环境空间缩为低维流形。只能写成速度级等式 $\mathbf A(\mathbf q)\dot{\mathbf q}=\mathbf 0$ 且\emph{不可积}（不存在 $\boldsymbol\phi(\mathbf q)$ 使 $\nabla\boldsymbol\phi\cdot\dot{\mathbf q}=\mathbf 0$ 与之等价）的约束称为\emph{非完整}（nonholonomic）约束。
\end{definition}

完整约束的典型是闭链机构的回路约束 $\boldsymbol\phi(\mathbf q)=\mathbf p_{\mathrm{end}_1}(\mathbf q)-\mathbf p_{\mathrm{end}_2}(\mathbf q)=\mathbf 0$；用最小坐标集时这类约束被坐标选择隐式满足，而闭链与接触需用 Lagrange 乘子显式处理（主体留 \cref{ch:constrained-dynamics}）。非完整约束的经典例是平面无滑滚动的差速轮：位形 $\mathbf q=(x,y,\theta)$，约束 $\dot x\sin\theta-\dot y\cos\theta=0$（不能瞬时侧移）不可积——机器人虽不能侧滑，却可经「先转再走」到达平面上任意点，这正是非完整系统强可达性的数学根源。非完整约束无法靠减少坐标消去，须以乘子并入
\begin{equation}
  \frac{\mathrm d}{\mathrm dt}\frac{\partial L}{\partial\dot{\mathbf q}}-\frac{\partial L}{\partial\mathbf q}=\boldsymbol\tau+\mathbf A^\top(\mathbf q)\boldsymbol\lambda.
  \label{eq:gm-nonholonomic}
\end{equation}

% ════════════════════════════════════════════════════════════════
\section{d'Alembert 原理与 Euler--Lagrange 方程}
\label{sec:gm-el}

\paragraph{动机：把 Newton 定律投影到约束允许的方向。}
直接对每个质点写 $\mathbf F_i=m_i\ddot{\mathbf r}_i$ 必然牵涉未知的关节约束力。d'Alembert 原理用\emph{虚位移}——给定时刻 $t$、与约束兼容的假想无穷小位移 $\delta\mathbf r_i$（非真实运动，是时间「冻结」时的测试方向）——把动力学问题转化为投影问题，让约束力自动消失。

\begin{theorem}[d'Alembert 原理与广义力]\label{thm:gm-dalembert}
设系统受\emph{理想约束}（约束力不做虚功，$\sum_i\mathbf R_i\cdot\delta\mathbf r_i=0$）。则对一切与约束兼容的虚位移，
\begin{equation}
  \sum_i\big(\mathbf F_i-m_i\ddot{\mathbf r}_i\big)\cdot\delta\mathbf r_i=0.
  \label{eq:gm-dalembert}
\end{equation}
以 $\mathbf r_i=\mathbf r_i(\mathbf q,t)$、$\delta\mathbf r_i=\sum_j(\partial\mathbf r_i/\partial q_j)\delta q_j$ 代入并用 $\delta q_j$ 独立，定义\emph{广义力} $Q_j=\sum_i\mathbf F_i\cdot\partial\mathbf r_i/\partial q_j$，即把 $3N$ 维 Newton 方程化为 $n$ 维广义坐标方程。
\end{theorem}

\begin{proof}
由 Newton 第二定律 $\mathbf F_i+\mathbf R_i=m_i\ddot{\mathbf r}_i$ 得 $\mathbf F_i-m_i\ddot{\mathbf r}_i=-\mathbf R_i$，故 $\sum_i(\mathbf F_i-m_i\ddot{\mathbf r}_i)\cdot\delta\mathbf r_i=-\sum_i\mathbf R_i\cdot\delta\mathbf r_i=0$，最后一步用理想约束假设。\cref{eq:gm-dalembert} 的威力在于约束力 $\mathbf R_i$ 从方程中彻底消失——对串联机器人，关节反力不出现在最终方程里。
\end{proof}

\paragraph{从虚功到 Hamilton 原理。}
要把惯性项 $\sum_i m_i\ddot{\mathbf r}_i\cdot\partial\mathbf r_i/\partial q_j$ 表为 $\mathbf q,\dot{\mathbf q},\ddot{\mathbf q}$ 的函数，最简洁的途径是变分原理。\cref{ch:opt-variational-pmp} 已把一般变分法到 Pontryagin 原理的整条逻辑链严格化，本节只取其针对机器人的特化。

\begin{theorem}[Hamilton 原理与 Euler--Lagrange 方程]\label{thm:gm-el}
定义 Lagrangian $L(\mathbf q,\dot{\mathbf q},t)=T-V$（动能减势能），作用量 $S=\int_{t_1}^{t_2}L\,\mathrm dt$。在所有连接固定端点 $\mathbf q(t_1),\mathbf q(t_2)$ 的路径中，真实路径使 $S$ 取驻值（$\delta S=0$），等价于
\begin{equation}
  \frac{\mathrm d}{\mathrm dt}\frac{\partial L}{\partial\dot q_i}-\frac{\partial L}{\partial q_i}=\tau_i,\qquad i=1,\dots,n,
  \label{eq:gm-euler-lagrange}
\end{equation}
其中 $\tau_i$ 为不能由势能导出的非保守广义力（旋转关节为力矩、平移关节为力）。
\end{theorem}

\begin{proof}
取扰动 $\mathbf q_\varepsilon=\mathbf q+\varepsilon\boldsymbol\eta$，$\boldsymbol\eta(t_1)=\boldsymbol\eta(t_2)=\mathbf 0$。一阶展开
$\delta S=\int_{t_1}^{t_2}\big[\tfrac{\partial L}{\partial q_i}\eta_i+\tfrac{\partial L}{\partial\dot q_i}\dot\eta_i\big]\mathrm dt$；对第二项分部积分，
$\int\tfrac{\partial L}{\partial\dot q_i}\dot\eta_i\,\mathrm dt=[\tfrac{\partial L}{\partial\dot q_i}\eta_i]_{t_1}^{t_2}-\int\tfrac{\mathrm d}{\mathrm dt}\tfrac{\partial L}{\partial\dot q_i}\eta_i\,\mathrm dt$，
边界项因端点固定而消失。于是 $\delta S=\int\big[\tfrac{\partial L}{\partial q_i}-\tfrac{\mathrm d}{\mathrm dt}\tfrac{\partial L}{\partial\dot q_i}\big]\eta_i\,\mathrm dt=0$ 对一切零端点 $\boldsymbol\eta$ 成立，由变分基本引理（\cref{lem:pmp-dubois}）得保守情形的 \cref{eq:gm-euler-lagrange}（$\tau_i=0$）。非保守广义力由 \cref{thm:gm-dalembert} 的 $Q_j$ 直接补在右端。
\end{proof}

\begin{insight}
为何 $L=T-V$ 而非 $T+V$？真实路径是「惯性」（想继续运动、$T$ 大）与「保守力」（想回到低势能、$V$ 小）之间的平衡，$T-V$ 的驻值条件恰表达这一平衡。而 Lagrange 法相对 Newton--Euler 的核心优势，正是 \cref{thm:gm-dalembert} 揭示的：理想约束力不做虚功，故\emph{无须计算关节间内力}——这与电路分析中网孔法自动消去支路电流如出一辙：选对「坐标」，一部分变量自动消失。
\end{insight}

% ════════════════════════════════════════════════════════════════
\section{质量矩阵 \texorpdfstring{$\mathbf M(\mathbf q)$}{M(q)} 及其性质}
\label{sec:gm-massmatrix}

\paragraph{两条构造路径。}
机器人动能恒为广义速度的二次型 $T=\tfrac12\dot{\mathbf q}^\top\mathbf M(\mathbf q)\dot{\mathbf q}$，$\mathbf M$ 即\emph{质量矩阵}。其一为 Jacobian 叠加法：第 $i$ 连杆质心速度 $\mathbf v_{c_i}=\mathbf J_{v_i}\dot{\mathbf q}$、角速度 $\boldsymbol\omega_i=\mathbf J_{\omega_i}\dot{\mathbf q}$，逐连杆动能求和得
\begin{equation}
  \mathbf M(\mathbf q)=\sum_{i=1}^n\Big(m_i\,\mathbf J_{v_i}^\top\mathbf J_{v_i}+\mathbf J_{\omega_i}^\top\,{}^0\mathbf I_{c_i}\,\mathbf J_{\omega_i}\Big),\qquad {}^0\mathbf I_{c_i}=\mathbf R_i\mathbf I_{c_i}^{b}\mathbf R_i^\top.
  \label{eq:gm-massmatrix}
\end{equation}
其二为空间向量法：以第 $i$ 连杆体速度 $\mathbf v_i\in\mathbb R^6$ 与 $6\times6$ 空间惯性 $\boldsymbol{\mathcal I}_i$ 写 $T=\tfrac12\sum_i\mathbf v_i^\top\boldsymbol{\mathcal I}_i\mathbf v_i$，沿关节约束展开后同样得 $T=\tfrac12\dot{\mathbf q}^\top\mathbf M\dot{\mathbf q}$，且 $\mathbf M$ 可由复合刚体算法（CRBA，\cref{ch:spatial-dynamics}）以 $O(n^2)$ 递推。前者宜手推（2--3 自由度），后者宜高效数值（任意自由度），二者给出同一 $\mathbf M$。

\begin{theorem}[$\mathbf M(\mathbf q)$ 对称正定]\label{thm:gm-spd}
对一切 $\mathbf q\in\mathcal Q$，串联机械臂的质量矩阵 $\mathbf M(\mathbf q)$ 对称正定。
\end{theorem}

\begin{proof}
对称性：\cref{eq:gm-massmatrix} 每项形如 $\mathbf A^\top\mathbf B\mathbf A$（$\mathbf B$ 对称），故和对称。正定性：对任意 $\dot{\mathbf q}\ne\mathbf0$，
$\dot{\mathbf q}^\top\mathbf M\dot{\mathbf q}=2T=\sum_i\big(m_i\|\mathbf J_{v_i}\dot{\mathbf q}\|^2+(\mathbf J_{\omega_i}\dot{\mathbf q})^\top\mathbf I_{c_i}(\mathbf J_{\omega_i}\dot{\mathbf q})\big)\ge0$，各项非负。要严格为正，注意串联结构下不存在 $\dot{\mathbf q}\ne\mathbf0$ 使所有连杆同时静止（第 $i$ 关节直接驱动连杆 $i$，对应 $\mathbf J_{\omega_i}$ 或 $\mathbf J_{v_i}$ 的第 $i$ 列非零）；更严格地，$\mathbf M\succ0$ 等价于堆叠 Jacobian $[\mathbf J_{v_1}^\top,\mathbf J_{\omega_1}^\top,\dots]^\top$ 列满秩，对串联臂在所有构型成立。
\end{proof}

\begin{insight}
$\mathbf M\succ0$ 的物理含义是「不存在零能量的运动」：只要关节有非零速度，系统就有正动能。这看似显然却对数值算法至关重要——它保证 $\mathbf M$ 可逆（正动力学 $\ddot{\mathbf q}=\mathbf M^{-1}(\cdots)$ 总有唯一解）、Cholesky 分解总成功。失效仅源于非物理输入或冗余约束：URDF 惯性参数取零或负使 $\mathbf M$ 退化；并联机构的冗余约束须先约去冗余坐标。\textbf{需厘清}：运动学奇异（Jacobian $\mathbf J$ 失秩）只使\emph{操作空间}惯性 $\boldsymbol\Lambda=(\mathbf J\mathbf M^{-1}\mathbf J^\top)^{-1}$ 退化，关节空间 $\mathbf M$ 本身仍正定（操作空间动力学主体见 \cref{ch:operational-space}\cite{khatib1987operational}）。
\end{insight}

% ════════════════════════════════════════════════════════════════
\section{Coriolis 矩阵、Christoffel 符号与 \texorpdfstring{$\dot{\mathbf M}-2\mathbf C$}{Mdot-2C} 反对称性}
\label{sec:gm-christoffel}

\paragraph{从 Euler--Lagrange 到标准形。}
对 $L=\tfrac12\dot{\mathbf q}^\top\mathbf M\dot{\mathbf q}-V$，逐项计算 \cref{eq:gm-euler-lagrange}：
$\tfrac{\mathrm d}{\mathrm dt}\tfrac{\partial L}{\partial\dot q_i}=\sum_j m_{ij}\ddot q_j+\sum_{j,k}\tfrac{\partial m_{ij}}{\partial q_k}\dot q_k\dot q_j$，
$\tfrac{\partial L}{\partial q_i}=\tfrac12\sum_{j,k}\tfrac{\partial m_{jk}}{\partial q_i}\dot q_j\dot q_k-\tfrac{\partial V}{\partial q_i}$。
对 $\sum_{j,k}\tfrac{\partial m_{ij}}{\partial q_k}\dot q_j\dot q_k$ 用哑指标对称化（交换 $j\leftrightarrow k$ 取平均），再与 $-\tfrac12\sum\tfrac{\partial m_{jk}}{\partial q_i}\dot q_j\dot q_k$ 合并，得二次项系数即第一类 Christoffel 符号。

\begin{definition}[Christoffel 符号与 Coriolis 矩阵]\label{def:gm-christoffel}
\emph{第一类 Christoffel 符号}与由其构造的 Coriolis 矩阵元为
\begin{equation}
  c_{ijk}(\mathbf q)=\frac12\Big(\frac{\partial m_{ij}}{\partial q_k}+\frac{\partial m_{ik}}{\partial q_j}-\frac{\partial m_{jk}}{\partial q_i}\Big),\qquad
  C_{ij}(\mathbf q,\dot{\mathbf q})=\sum_k c_{ijk}(\mathbf q)\,\dot q_k.
  \label{eq:gm-christoffel}
\end{equation}
\end{definition}

\begin{theorem}[Christoffel 即位形空间的 Levi--Civita 联络]\label{thm:gm-levicivita}
将位形空间 $\mathcal Q$ 赋以 Riemann 度量 $g_{ij}=m_{ij}$（以质量矩阵为度量张量），则 \cref{eq:gm-christoffel} 恰是该度量的 Levi--Civita 联络系数；无外力时（$\boldsymbol\tau=\mathbf0,\,V=\mathrm{const}$）运动方程 $\mathbf M\ddot{\mathbf q}+\mathbf C\dot{\mathbf q}=\mathbf0$ 即\emph{测地线方程}——自由运动沿位形空间的「直线」行进。
\end{theorem}

\begin{insight}
这绝非记号巧合。Coriolis 力与离心力不是「真实的力」，而是因广义坐标系弯曲／旋转而生的表观力，故可完全用度量张量的偏导数表达——其角色与广义相对论中 Christoffel 符号描述「坐标系弯曲导致的表观力」完全一致（\cref{ch:matderiv} 给出度量与联络的微分几何背景）。这一视角将在 \cref{sec:gm-euler-poincare} 升级：李群上 Christoffel 项被余伴随作用 $\operatorname{ad}^*_{\boldsymbol\xi}$ 取代，二者是同一几何效应的坐标版与坐标无关版。
\end{insight}

\begin{theorem}[$\dot{\mathbf M}-2\mathbf C$ 反对称性]\label{thm:gm-skew}
记 $\mathbf N(\mathbf q,\dot{\mathbf q}):=\dot{\mathbf M}-2\mathbf C$。则有两个层次的结论：
\begin{description}
  \item[强版本] 当 $\mathbf C$ 由 \cref{eq:gm-christoffel} 的 Christoffel 构造时，$\mathbf N$ 反对称：$\mathbf N^\top=-\mathbf N$，即 $\mathbf x^\top\mathbf N\mathbf x=0$ 对\emph{一切} $\mathbf x\in\mathbb R^n$。
  \item[弱版本] 对任意使 $\mathbf C\dot{\mathbf q}$ 正确的合法 $\mathbf C$，恒有 $\dot{\mathbf q}^\top\mathbf N\dot{\mathbf q}=0$（能量守恒形式），但 $\mathbf N$ 未必反对称。
\end{description}
\end{theorem}

\begin{proof}
计算 $\dot m_{ij}=\sum_k\tfrac{\partial m_{ij}}{\partial q_k}\dot q_k$，故
\[
N_{ij}=\dot m_{ij}-2C_{ij}=\sum_k\Big[\tfrac{\partial m_{ij}}{\partial q_k}-\big(\tfrac{\partial m_{ij}}{\partial q_k}+\tfrac{\partial m_{ik}}{\partial q_j}-\tfrac{\partial m_{jk}}{\partial q_i}\big)\Big]\dot q_k=\sum_k\Big(\tfrac{\partial m_{jk}}{\partial q_i}-\tfrac{\partial m_{ik}}{\partial q_j}\Big)\dot q_k.
\]
交换指标 $i\leftrightarrow j$ 即得 $N_{ji}=-N_{ij}$，强版本成立。弱版本是其即时推论（反对称矩阵的二次型为零），且对非 Christoffel 的 $\mathbf C$ 仍成立，因 $\dot{\mathbf q}^\top\mathbf N\dot{\mathbf q}=\tfrac{\mathrm d}{\mathrm dt}T-\dot{\mathbf q}^\top\mathbf C\dot{\mathbf q}$ 是能量平衡的恒等式，与 $\mathbf C$ 的具体选取无关。
\end{proof}

\begin{pitfall}[误区：「任何合法的 $\mathbf C$ 都使 $\dot{\mathbf M}-2\mathbf C$ 反对称」]
$\mathbf C$ 并不唯一：给定 $\mathbf M$，只要 $\mathbf C\dot{\mathbf q}$ 等于正确的 Coriolis/离心项，$\mathbf C$ 即合法——例如 $\mathbf C'=\mathbf C+\mathbf S$（任意满足 $\mathbf S\dot{\mathbf q}=\mathbf0$ 的 $\mathbf S$）仍合法。但只有 Christoffel 构造保证\emph{全向量}反对称（强版本）。这一区分有实际后果：Slotine--Li 自适应控制\cite{slotine1987adaptive} 的 Lyapunov 证明需要 $\mathbf s^\top\mathbf N\mathbf s=0$（其中合成误差 $\mathbf s\ne\dot{\mathbf q}$），只有强版本能提供；用非 Christoffel 的 $\mathbf C$ 则证明失效。凡需反对称性的控制器设计，一律采用 Christoffel 构造的 $\mathbf C$。
\end{pitfall}

\begin{remark}[同一 Coriolis 的三种视角]
本章 Christoffel 解析视角（\cref{def:gm-christoffel}）、\cref{ch:spatial-dynamics} 的递推视角（Coriolis 藏在力叉积 $\mathbf v\times^*(\boldsymbol{\mathcal I}\mathbf v)$ 中）、以及 \cref{sec:gm-euler-poincare} 的李群视角（$\operatorname{ad}^*_{\boldsymbol\xi}$ 项）描述的是同一个物理对象。三者互引：$c_{ijk}$ 是 $n$ 维关节空间的展开，$\mathbf v\times^*(\boldsymbol{\mathcal I}\mathbf v)$ 是 $O(n)$ 递推的体现，$\operatorname{ad}^*_{\boldsymbol\xi}\boldsymbol\mu$ 是李群约化下的坐标无关形式。
\end{remark}

% ════════════════════════════════════════════════════════════════
\section{重力项、被动性接口与 2R 算例}
\label{sec:gm-gravity}

\paragraph{重力项与被动性。}
势能为各连杆重力势之和 $V(\mathbf q)=\sum_i m_i\mathbf g_0^\top\mathbf p_{c_i}(\mathbf q)$（$\mathbf g_0$ 为重力加速度向量），重力项 $\mathbf g(\mathbf q)=\partial V/\partial\mathbf q$。$\dot{\mathbf M}-2\mathbf C$ 反对称的最直接推论是\emph{被动性}：

\begin{theorem}[被动性]\label{thm:gm-passivity}
取储能函数 $H=\tfrac12\dot{\mathbf q}^\top\mathbf M\dot{\mathbf q}+V$，则沿 \cref{eq:gm-standard} 有 $\dot H=\dot{\mathbf q}^\top\boldsymbol\tau$，即系统以 $\dot{\mathbf q}$ 为输出、$\boldsymbol\tau$ 为输入是无源的。
\end{theorem}

\begin{proof}
$\dot H=\dot{\mathbf q}^\top\mathbf M\ddot{\mathbf q}+\tfrac12\dot{\mathbf q}^\top\dot{\mathbf M}\dot{\mathbf q}+\dot{\mathbf q}^\top\mathbf g$。由 \cref{eq:gm-standard} 代 $\mathbf M\ddot{\mathbf q}=\boldsymbol\tau-\mathbf C\dot{\mathbf q}-\mathbf g$，得 $\dot H=\dot{\mathbf q}^\top\boldsymbol\tau+\tfrac12\dot{\mathbf q}^\top(\dot{\mathbf M}-2\mathbf C)\dot{\mathbf q}$，第二项由 \cref{thm:gm-skew} 弱版本为零。
\end{proof}

\paragraph{参数线性性与接口陈述。}
\cref{eq:gm-standard} 的左端对惯性参数 $\boldsymbol\theta$ 线性：$\mathbf M\ddot{\mathbf q}+\mathbf C\dot{\mathbf q}+\mathbf g=\mathbf Y(\mathbf q,\dot{\mathbf q},\ddot{\mathbf q})\boldsymbol\theta$，回归量 $\mathbf Y$ 不含 $\boldsymbol\theta$。此结构使自适应控制与系统辨识可行——这是\emph{结构性接口}，控制器设计与 Lyapunov 证明主体留 \cref{ch:lyapunov-stability}、\cref{ch:operational-space}，本章只陈述结构本身：被动性（\cref{thm:gm-passivity}）$\Rightarrow$ 无源控制；$\dot{\mathbf M}-2\mathbf C$ 强版本 $\Rightarrow$ Slotine--Li 自适应；参数线性 $\mathbf Y\boldsymbol\theta\Rightarrow$ 辨识可行。

\begin{example}[2R 平面臂：闭式 $\mathbf M,\mathbf C,\mathbf g$ 与反对称验证]\label{ex:gm-2r}
两旋转关节在竖直平面，连杆长 $\ell_1,\ell_2$、末端集中质量 $m_1,m_2$，关节角 $q_1$（对水平）、$q_2$（相对连杆 1）。记 $c_2=\cos q_2,\,s_2=\sin q_2,\,h=-m_2\ell_1\ell_2 s_2$。逐质心位置求导得动能，整理为 $T=\tfrac12\dot{\mathbf q}^\top\mathbf M\dot{\mathbf q}$：
\begin{equation}
  \mathbf M=\begin{pmatrix}(m_1+m_2)\ell_1^2+m_2\ell_2^2+2m_2\ell_1\ell_2 c_2 & m_2\ell_2^2+m_2\ell_1\ell_2 c_2\\[2pt] m_2\ell_2^2+m_2\ell_1\ell_2 c_2 & m_2\ell_2^2\end{pmatrix}.
  \label{eq:gm-2r-M}
\end{equation}
由 \cref{eq:gm-christoffel}（仅 $\partial m_{ij}/\partial q_2$ 非零），Coriolis 与重力为
\begin{equation}
  \mathbf C=\begin{pmatrix}h\dot q_2 & h(\dot q_1+\dot q_2)\\ -h\dot q_1 & 0\end{pmatrix},\quad
  \mathbf g=\begin{pmatrix}(m_1+m_2)g\ell_1\cos q_1+m_2 g\ell_2\cos(q_1+q_2)\\ m_2 g\ell_2\cos(q_1+q_2)\end{pmatrix}.
  \label{eq:gm-2r-Cg}
\end{equation}
验证：$\dot{\mathbf M}=\bigl(\begin{smallmatrix}2h\dot q_2 & h\dot q_2\\ h\dot q_2 & 0\end{smallmatrix}\bigr)$，故
$\dot{\mathbf M}-2\mathbf C=\bigl(\begin{smallmatrix}0 & -h(2\dot q_1+\dot q_2)\\ h(2\dot q_1+\dot q_2) & 0\end{smallmatrix}\bigr)$，确为反对称（$N_{12}=-N_{21},\,N_{11}=N_{22}=0$），印证 \cref{thm:gm-skew} 强版本。
\end{example}

\begin{remark}[为何 $\ge3$ 自由度不手推]
2R 闭式尚可手算，但 3R 空间臂的 $\mathbf M$ 含数十个三角乘积项，Christoffel 符号 $c_{ijk}$ 有 $n^3$ 个分量，逐项展开呈组合爆炸。工程上 $\ge3$ 自由度一律用符号计算工具或 \cref{ch:spatial-dynamics} 的 $O(n)$ 递推数值生成，手推仅用于建立 2R 这样的最小直觉。
\end{remark}

% ════════════════════════════════════════════════════════════════
% 第二部分 — 几何力学
% ════════════════════════════════════════════════════════════════
\section{为何需要几何力学：浮基系统的拓扑硬约束}
\label{sec:gm-why-geom}

\paragraph{动机：基座的「位置」无法无奇异地参数化。}
\crefrange{sec:gm-history}{sec:gm-gravity} 的 Euler--Lagrange 框架默认 $\mathcal Q$ 局部同胚于 $\mathbb R^n$、$\dot{\mathbf q}$ 与广义力同维。对浮基系统（四足、人形、自由漂浮空间臂）这一前提失效：其位形空间天然是 $\mathcal Q=\SE(3)\times\mathbb T^n$，而 $\SE(3)$ 中的 $\SO(3)$ 因子无法被 $\mathbb R^3$ 无奇异地全覆盖（毛球定理的推论）。具体困境有三：

\begin{enumerate}
  \item \textbf{欧拉角奇异（万向锁）}：ZYX 欧拉角 $(\phi,\theta,\psi)$ 与体角速度的映射在 $\theta=\pi/2$ 处的矩阵奇异——物理角速度有限，欧拉角速率却趋于无穷，$\mathbf M$ 与 Christoffel 符号在此奇异，数值积分崩溃。
  \item \textbf{伪张量}：以欧拉角写出的 $\mathbf M$ 含大量 $\sin/\cos$ 项，Christoffel 系数在不同欧拉角约定下形式迥异——换一种约定即须重推。
  \item \textbf{维度不匹配}：用单位四元数避开奇异，则位形维 $n_q=7+n$ 但速度维 $n_v=6+n$（四元数 4 参数 $+1$ 约束 $\|\mathbf q\|=1$），标准 Euler--Lagrange 因 $\dot{\mathbf q}$ 与 $\boldsymbol\tau$ 维度不符而直接失效，问题沦为微分代数方程（DAE）。
\end{enumerate}

\paragraph{几何力学的承诺。}
Euler--Poincaré 方程直接在体坐标系工作，用李代数元 $\boldsymbol\xi=\mathbf g^{-1}\dot{\mathbf g}\in\se(3)$ 代替 $\dot{\mathbf q}$、用余伴随作用 $\operatorname{ad}^*_{\boldsymbol\xi}$ 代替 Christoffel 项，因\emph{根本不使用任何参数化}而天然避免表示奇异。其实际工程后果（详见后文）包括：InEKF 一致性需要 $\SE(3)$ 上的不变动力学作为过程模型\cite{barrau2017invariant}；质心动力学本质是 $\SE(3)$ 上的约化\cite{orin2013centroidal}；\cref{ch:spatial-dynamics} 的空间向量代数正是 $\se(3)/\se(3)^*$ 的矩阵实现。可类比地球表面运动：经纬度在极点失效，而「脚下局部切平面」永远有效——Euler--Poincaré 即此方案的力学版，始终在体坐标系的切空间（李代数）工作。

% ════════════════════════════════════════════════════════════════
\section{体速度、空间速度与左不变约化}
\label{sec:gm-bodyvel}

\paragraph{两种速度的精确定义。}
设系统位形为李群曲线 $\mathbf g(t)\in G$。\emph{体速度}（左平凡化）$\boldsymbol\xi=\mathbf g^{-1}\dot{\mathbf g}\in\mathfrak g$ 是「站在刚体上感受到的速度」；\emph{空间速度}（右平凡化）$\boldsymbol\eta=\dot{\mathbf g}\mathbf g^{-1}\in\mathfrak g$ 是「站在固定参考系看到的速度」。二者经伴随表示相联：
\begin{equation}
  \boldsymbol\eta=\dot{\mathbf g}\mathbf g^{-1}=\mathbf g(\mathbf g^{-1}\dot{\mathbf g})\mathbf g^{-1}=\operatorname{Ad}_{\mathbf g}\boldsymbol\xi,
  \label{eq:gm-ad-eta-xi}
\end{equation}
其中 $\operatorname{Ad}_{\mathbf g}$ 取自 \cref{ch:lie}。对 $\SE(3)$，$\mathbf g=(\mathbf R,\mathbf p)$ 的齐次矩阵给出
$\boldsymbol\omega_b=(\mathbf R^\top\dot{\mathbf R})^\vee,\,\mathbf v_b=\mathbf R^\top\dot{\mathbf p}$（体速度），而空间速度的线部分 $\mathbf v_s=\dot{\mathbf p}-\boldsymbol\omega_s\times\mathbf p$ 是「空间参考点（原点固定）处的速度」。

\begin{pitfall}[$\mathbf v_s$ 不是质心线速度的空间表达]
体速度的 $\mathbf v_b=\mathbf R^\top\dot{\mathbf p}$ 确是质心线速度在体坐标系的投影；但空间速度的 $\mathbf v_s=\dot{\mathbf p}-\boldsymbol\omega_s\times\mathbf p$ \emph{不是}质心线速度在世界系的表达，而是「原点固定处的速度」。这是 $6$D 空间速度最易混淆处——它是 $\se(3)^*$ 对偶配对意义下、计算功率时与 wrench 对应的对象，而非朴素的「质心速度向量」。
\end{pitfall}

\paragraph{左不变约化。}
刚体动能只依赖角速度与质心线速度，不依赖刚体在空间的具体位姿——即 $L(\mathbf g,\dot{\mathbf g})$ 在 $G$ 的左平移下不变：$L(\mathbf h\mathbf g,\mathbf h\dot{\mathbf g})=L(\mathbf g,\dot{\mathbf g})$。于是 $L$ 只依赖体速度 $\boldsymbol\xi$，\emph{下降}到李代数上的\emph{约化 Lagrangian}
\begin{equation}
  \ell:\mathfrak g\to\mathbb R,\qquad \ell(\boldsymbol\xi):=L(\mathbf e,\boldsymbol\xi).
  \label{eq:gm-reduced-lagrangian}
\end{equation}
对 $\SE(3)$，原 $L$ 定义在切丛 $TG$（12 维）上，约化后 $\ell$ 只在 $\mathfrak g$（6 维）上——这就是「利用对称性降维」。

\begin{table}[H]\centering\small
\caption{两套 $6$D 分量序对照（本章 $\leftrightarrow$ 本书 $\se(3)$）：由置换 $\mathbf P=\bigl(\begin{smallmatrix}\mathbf0&\mathbf I\\\mathbf I&\mathbf0\end{smallmatrix}\bigr)$ 互换}
\label{tab:gm-ordering}
\begin{tabular}{lll}
\toprule
 & 本章几何力学（角在前） & 本书 \cref{ch:lie} 之 $\se(3)$（线在前）\\
\midrule
速度坐标 & $\boldsymbol\xi=[\boldsymbol\omega;\mathbf v]$ & $\boldsymbol\xi_{\mathrm{book}}=[\boldsymbol\rho;\boldsymbol\phi]$（坐标率）\\
动量坐标 & $\boldsymbol\mu=[\mathbf m;\mathbf f]$（角动量；线动量） & 对偶同序互换\\
$\operatorname{ad}_{\boldsymbol\xi}$ 块 & 见 \cref{eq:gm-ad-se3}（下三角块为 $[\mathbf v]_\times$）& $\mathbf P\operatorname{ad}_{\boldsymbol\xi}\mathbf P$（上三角块）\\
来源 & Featherstone, Lynch--Park, Marsden--Ratiu & Sophus/十四讲/Barfoot 约定\\
\bottomrule
\end{tabular}
\end{table}

\begin{remark}[「left-invariant 约化」$\ne$「左扰动」]
\cref{ch:lie} 的扰动方向（本书主线\emph{右扰动} $\mathbf R\,\mathrm{Exp}(\delta\boldsymbol\phi)$）与本节「left-invariant reduction」（体速度 $\boldsymbol\xi=\mathbf g^{-1}\dot{\mathbf g}$）是\emph{正交}的两个概念，不可混淆：前者是「对群参数求导时把扰动加在哪一侧」，后者是「约化时把对称性除在哪一侧」。体坐标系约化是几何力学的标准选择，与本书右扰动求导约定并不冲突；凡涉及对位形 $\mathbf g$ 求导（如指向 \cref{ch:constrained-dynamics} 解析微分的钩子）仍采本书右扰动、右雅可比 $\mathbf J_r$。
\end{remark}

% ════════════════════════════════════════════════════════════════
\section{约束变分与 Euler--Poincaré 方程}
\label{sec:gm-euler-poincare}

\paragraph{动机：约化后的变分被群结构锁死。}
标准 Hamilton 原理中变分 $\delta\mathbf q$ 是自由的（端点为零即可）。但约化到李代数后，速度变分 $\delta\boldsymbol\xi$ \emph{不自由}——它被一个约束方程锁定。理解这一约束是推导 Euler--Poincaré 方程最精妙的一步。

\begin{lemma}[约束变分]\label{lem:gm-constrained-variation}
取两参数族 $\mathbf g(t,\varepsilon)$，定义 $\boldsymbol\xi=\mathbf g^{-1}\partial_t\mathbf g$、$\boldsymbol\eta=\mathbf g^{-1}\partial_\varepsilon\mathbf g$（均 $\in\mathfrak g$）。则
\begin{equation}
  \delta\boldsymbol\xi=\dot{\boldsymbol\eta}+[\boldsymbol\xi,\boldsymbol\eta]=\dot{\boldsymbol\eta}+\operatorname{ad}_{\boldsymbol\xi}\boldsymbol\eta,\qquad \delta\boldsymbol\xi:=\partial_\varepsilon\boldsymbol\xi|_{\varepsilon=0}.
  \label{eq:gm-constrained-variation}
\end{equation}
\end{lemma}

\begin{proof}
对矩阵李群直接验证。由 $\mathbf g\mathbf g^{-1}=\mathbf I$ 对 $\varepsilon$ 求导得 $\partial_\varepsilon\mathbf g^{-1}=-\mathbf g^{-1}(\partial_\varepsilon\mathbf g)\mathbf g^{-1}=-\boldsymbol\eta\mathbf g^{-1}$。于是
$\delta\boldsymbol\xi=\partial_\varepsilon(\mathbf g^{-1}\partial_t\mathbf g)=(\partial_\varepsilon\mathbf g^{-1})\partial_t\mathbf g+\mathbf g^{-1}\partial_{\varepsilon t}^2\mathbf g$。
第一项 $=-\boldsymbol\eta\mathbf g^{-1}\partial_t\mathbf g=-\boldsymbol\eta\boldsymbol\xi$。第二项用混合偏导交换 $\partial_{\varepsilon t}\mathbf g=\partial_{t\varepsilon}\mathbf g$：
$\mathbf g^{-1}\partial_t(\partial_\varepsilon\mathbf g)=\mathbf g^{-1}\partial_t(\mathbf g\boldsymbol\eta)=\mathbf g^{-1}(\dot{\mathbf g}\boldsymbol\eta+\mathbf g\dot{\boldsymbol\eta})=\boldsymbol\xi\boldsymbol\eta+\dot{\boldsymbol\eta}$。
合并得 $\delta\boldsymbol\xi=-\boldsymbol\eta\boldsymbol\xi+\boldsymbol\xi\boldsymbol\eta+\dot{\boldsymbol\eta}=[\boldsymbol\xi,\boldsymbol\eta]+\dot{\boldsymbol\eta}$。
\end{proof}

\begin{insight}
约束变分的物理意义：变分方向 $\boldsymbol\eta$ 可任选（端点为零），但由它引起的速度变分 $\delta\boldsymbol\xi$ 被群结构通过 \cref{eq:gm-constrained-variation} 完全锁死——如同不可伸长的绳子把两质点的运动绑在一起，这里「绳子」就是群结构本身。若 $\delta\boldsymbol\xi$ 像 $\mathbb R^n$ 中那样自由，Hamilton 原理只会给出动量守恒 $\tfrac{\mathrm d}{\mathrm dt}\partial\ell/\partial\boldsymbol\xi=\mathbf0$；这对自由刚体的\emph{空间}角动量成立，但对\emph{体}角动量不成立（其方向在变）。约束变分恰好补回差异，产生额外的 $\operatorname{ad}^*_{\boldsymbol\xi}$ 项，描述的正是「体坐标系本身在旋转」这一事实。$\SO(3)$ 特例 $[\boldsymbol\xi,\boldsymbol\eta]=\boldsymbol\omega\times\boldsymbol\eta$，故 $\delta\boldsymbol\omega=\dot{\boldsymbol\eta}+\boldsymbol\omega\times\boldsymbol\eta$。
\end{insight}

\begin{theorem}[Euler--Poincaré 方程]\label{thm:gm-euler-poincare}
对左不变 Lagrangian 的约化 $\ell(\boldsymbol\xi)$，Hamilton 原理 $\delta\int\ell\,\mathrm dt=0$ 配以约束变分 \cref{eq:gm-constrained-variation} 给出
\begin{equation}
  \boxed{\;\frac{\mathrm d}{\mathrm dt}\frac{\partial\ell}{\partial\boldsymbol\xi}=\operatorname{ad}^*_{\boldsymbol\xi}\frac{\partial\ell}{\partial\boldsymbol\xi}+\boldsymbol\tau\;},\qquad \boldsymbol\mu:=\frac{\partial\ell}{\partial\boldsymbol\xi}\in\mathfrak g^*,
  \label{eq:gm-euler-poincare}
\end{equation}
其中余伴随 $\operatorname{ad}^*_{\boldsymbol\xi}$ 由 $\langle\operatorname{ad}^*_{\boldsymbol\xi}\boldsymbol\mu,\boldsymbol\eta\rangle=\langle\boldsymbol\mu,\operatorname{ad}_{\boldsymbol\xi}\boldsymbol\eta\rangle$ 定义，$\boldsymbol\tau\in\mathfrak g^*$ 为外广义力。该方程源出 Poincaré（1901）\cite{poincare1901forme}。
\end{theorem}

\begin{proof}
$\delta S=\int\langle\partial\ell/\partial\boldsymbol\xi,\delta\boldsymbol\xi\rangle\,\mathrm dt=\int\langle\boldsymbol\mu,\dot{\boldsymbol\eta}+\operatorname{ad}_{\boldsymbol\xi}\boldsymbol\eta\rangle\,\mathrm dt$。对 $\langle\boldsymbol\mu,\dot{\boldsymbol\eta}\rangle$ 项分部积分，端点 $\boldsymbol\eta(t_0)=\boldsymbol\eta(t_1)=\mathbf0$ 使边界项消失，得 $-\int\langle\tfrac{\mathrm d}{\mathrm dt}\boldsymbol\mu,\boldsymbol\eta\rangle\,\mathrm dt$；对 $\langle\boldsymbol\mu,\operatorname{ad}_{\boldsymbol\xi}\boldsymbol\eta\rangle$ 用余伴随定义化为 $\langle\operatorname{ad}^*_{\boldsymbol\xi}\boldsymbol\mu,\boldsymbol\eta\rangle$。合并 $\delta S=\int\langle-\tfrac{\mathrm d}{\mathrm dt}\boldsymbol\mu+\operatorname{ad}^*_{\boldsymbol\xi}\boldsymbol\mu,\boldsymbol\eta\rangle\,\mathrm dt=0$，由 $\boldsymbol\eta$ 任意（端点零）得 \cref{eq:gm-euler-poincare}（外力 $\boldsymbol\tau$ 由非保守广义力补在右端）。
\end{proof}

\begin{insight}
\cref{eq:gm-euler-poincare} 各项的物理意义：$\boldsymbol\mu=\partial\ell/\partial\boldsymbol\xi$ 是体坐标系广义动量（对应 $\mathbf p=\partial L/\partial\dot{\mathbf q}$）；$\tfrac{\mathrm d}{\mathrm dt}\boldsymbol\mu$ 是其时间导数；\textbf{$\operatorname{ad}^*_{\boldsymbol\xi}\boldsymbol\mu$ 不是新物理，而是体坐标系旋转的几何效应——它与 $\mathbb R^n$ 中的 Christoffel 项扮演完全相同的角色}（\cref{thm:gm-levicivita}），只是在李群语言中表达得坐标无关。Euler--Poincaré 与标准 Euler--Lagrange 的本质区别不是「换记号」：后者活在切丛 $TG$（$2n$ 维），前者活在李代数 $\mathfrak g$（$n$ 维），维度直接减半；代价是从 $\boldsymbol\xi(t)$ 恢复位形 $\mathbf g(t)$ 需额外积分重构方程 $\dot{\mathbf g}=\mathbf g\boldsymbol\xi$。注意体形式（左约化）与空间形式（右约化，$\tfrac{\mathrm d}{\mathrm dt}\boldsymbol\nu=-\operatorname{ad}^*_{\boldsymbol\eta}\boldsymbol\nu+\boldsymbol\tau_s$）\emph{符号相反}：机器人学多用体形式（Featherstone/Pinocchio），流体力学的 Euler 方程是右约化结果。
\end{insight}

% ════════════════════════════════════════════════════════════════
\section{退化验证：SO(3) Euler 方程与 SE(3) 六维 Newton--Euler}
\label{sec:gm-degenerate}

\paragraph{SO(3)：Euler 刚体方程作为推论。}
在 $\SO(3)$ 上 $\mathfrak g\cong(\mathbb R^3,\times)$，$\operatorname{ad}_{\boldsymbol\xi}\boldsymbol\eta=\boldsymbol\xi\times\boldsymbol\eta$。由 $\langle\operatorname{ad}^*_{\boldsymbol\xi}\boldsymbol\mu,\boldsymbol\eta\rangle=\boldsymbol\mu\cdot(\boldsymbol\xi\times\boldsymbol\eta)=(\boldsymbol\mu\times\boldsymbol\xi)\cdot\boldsymbol\eta$ 得 $\operatorname{ad}^*_{\boldsymbol\omega}\boldsymbol\pi=\boldsymbol\pi\times\boldsymbol\omega$。取 $\ell(\boldsymbol\omega)=\tfrac12\boldsymbol\omega^\top\mathbf I_c\boldsymbol\omega$、$\boldsymbol\pi=\mathbf I_c\boldsymbol\omega$，代入 \cref{eq:gm-euler-poincare}：
\begin{equation}
  \frac{\mathrm d}{\mathrm dt}(\mathbf I_c\boldsymbol\omega)=(\mathbf I_c\boldsymbol\omega)\times\boldsymbol\omega+\boldsymbol\tau
  \;\Longleftrightarrow\;
  \boxed{\;\mathbf I_c\dot{\boldsymbol\omega}+\boldsymbol\omega\times(\mathbf I_c\boldsymbol\omega)=\boldsymbol\tau\;}.
  \label{eq:gm-euler-rigid}
\end{equation}
经典 Euler 刚体方程不是独立的物理定律，而是 $\SO(3)$ 上 Euler--Poincaré 方程的自动推论。自由旋转（$\boldsymbol\tau=\mathbf0$）时体角动量模长守恒（$\tfrac{\mathrm d}{\mathrm dt}|\boldsymbol\pi|^2=2\boldsymbol\pi\cdot(\boldsymbol\pi\times\mathbf I_c^{-1}\boldsymbol\pi)=0$）、动能守恒（$\tfrac{\mathrm d}{\mathrm dt}T=\boldsymbol\omega\cdot(\boldsymbol\pi\times\boldsymbol\omega)=0$），两守恒量把 3D 相空间约束到 1D 曲线——这预告了 \cref{sec:gm-liepoisson} 的 Poinsot 构造。

\begin{pitfall}[$\boldsymbol\omega\times(\mathbf I_c\boldsymbol\omega)$ 不是物理力矩]
方程 $\mathbf I_c\dot{\boldsymbol\omega}=\boldsymbol\tau-\boldsymbol\omega\times(\mathbf I_c\boldsymbol\omega)$ 中右端第二项常被误认为「某种 Coriolis 力矩」。实际上它不是真实的物理力矩——在\emph{空间}坐标系中自由刚体的方程是 $\tfrac{\mathrm d}{\mathrm dt}(\mathbf R\mathbf I_c\boldsymbol\omega)=\mathbf0$（空间角动量守恒，无额外项）；额外项纯粹来自「体坐标系本身在旋转」，正是 $\operatorname{ad}^*_{\boldsymbol\xi}\boldsymbol\mu$ 的几何含义。
\end{pitfall}

\paragraph{SE(3)：六维 Newton--Euler。}
取 $\mathfrak g=\se(3)\cong\mathbb R^6$、体速度 $\boldsymbol\xi=[\boldsymbol\omega;\mathbf v]$。$\se(3)$ 的 Lie bracket（Featherstone 的 motion cross）及其 $6\times6$ 矩阵为
\begin{equation}
  \operatorname{ad}_{\boldsymbol\xi}=\begin{pmatrix}[\boldsymbol\omega]_\times & \mathbf0\\ [\mathbf v]_\times & [\boldsymbol\omega]_\times\end{pmatrix},\qquad
  \operatorname{ad}^*_{\boldsymbol\xi}=\operatorname{ad}_{\boldsymbol\xi}^\top=\begin{pmatrix}[\boldsymbol\omega]_\times & [\mathbf v]_\times\\ \mathbf0 & [\boldsymbol\omega]_\times\end{pmatrix},
  \label{eq:gm-ad-se3}
\end{equation}
此处 $\operatorname{ad}^*=\operatorname{ad}^\top$ 采\textbf{数学约定}（与 \cref{ch:lie} 一致）。以体坐标系空间惯性 $\boldsymbol{\mathcal I}$ 写 $\ell(\boldsymbol\xi)=\tfrac12\boldsymbol\xi^\top\boldsymbol{\mathcal I}\boldsymbol\xi$、$\boldsymbol\mu=\boldsymbol{\mathcal I}\boldsymbol\xi$，代入 \cref{eq:gm-euler-poincare} 得 $\boldsymbol{\mathcal I}\dot{\boldsymbol\xi}=\operatorname{ad}^*_{\boldsymbol\xi}(\boldsymbol{\mathcal I}\boldsymbol\xi)+\boldsymbol{\mathcal F}$。质心处展开（$\boldsymbol\mu=(\mathbf I_c\boldsymbol\omega;m\mathbf v)$、$\mathbf v\times m\mathbf v=\mathbf0$）即
\begin{equation}
  \mathbf I_c\dot{\boldsymbol\omega}+\boldsymbol\omega\times(\mathbf I_c\boldsymbol\omega)=\boldsymbol\tau_{\mathrm{ext}},\qquad
  m\dot{\mathbf v}+\boldsymbol\omega\times(m\mathbf v)=\mathbf f_{\mathrm{ext}}.
  \label{eq:gm-se3-newton-euler}
\end{equation}

\begin{pitfall}[math $\operatorname{ad}^*$ vs Featherstone force cross：本章头号符号易错点]
Featherstone 的 force cross 定义为相反符号 $\boldsymbol\xi\times^*:=-\operatorname{ad}^*_{\boldsymbol\xi}=-\operatorname{ad}_{\boldsymbol\xi}^\top$。两者文献中均常见：Murray--Li--Sastry、Marsden--Ratiu、本章 Euler--Poincaré 用数学 $\operatorname{ad}^*=\operatorname{ad}^\top$；\cref{ch:spatial-dynamics} 的 RNEA 力回溯用 $\times^*=-\operatorname{ad}^\top$。\textbf{混用是最常见的符号错误}：若误以 $\times^*$ 代入 \cref{eq:gm-euler-poincare} 会得到错号的 $\mathbf I_c\dot{\boldsymbol\omega}-\boldsymbol\omega\times\mathbf I_c\boldsymbol\omega=\boldsymbol\tau$。用 Featherstone 记号重写 EP 时 force cross 移到等号左边：$\boldsymbol{\mathcal I}\dot{\boldsymbol\xi}+\boldsymbol\xi\times^*(\boldsymbol{\mathcal I}\boldsymbol\xi)=\boldsymbol{\mathcal F}_{\mathrm{ext}}$，与 \cref{ch:spatial-dynamics} 单刚体统一式一致，亦与 Lynch--Park (8.40) 的 $\boldsymbol{\mathcal I}\dot{\boldsymbol{\mathcal V}}_b-[\operatorname{ad}_{\boldsymbol{\mathcal V}_b}]^\top\boldsymbol{\mathcal I}\boldsymbol{\mathcal V}_b=\boldsymbol{\mathcal F}_b$ 对应\cite{lynch2017modern}。本章每个 $\operatorname{ad}^*$ 公式均注明用的是哪套约定。
\end{pitfall}

\begin{insight}
整个 $6$D Newton--Euler 方程就是 $\SE(3)$ 上的 Euler--Poincaré。Featherstone 未用李群语言，但其空间向量代数编码了完全相同的几何内容：motion cross $=\operatorname{ad}$，force cross $\times^*=-\operatorname{ad}^\top$，spatial inertia $=$ 从 $\se(3)$ 到 $\se(3)^*$ 的映射。\cref{ch:spatial-dynamics}（计算视角）给出空间向量展开，本章（分析力学视角）给出其 Euler--Poincaré 由来，二者是同一对象的两面。
\end{insight}

% ════════════════════════════════════════════════════════════════
\section{浮基多体动力学}
\label{sec:gm-floating-base}

\paragraph{混合形式方程。}
浮基系统位形 $\mathcal Q=\SE(3)\times\mathbb T^n$，广义速度 $\mathbf v=(\boldsymbol\xi_b;\dot{\mathbf q}_j)\in\mathbb R^{6+n}$（基座体速度 $6$ 维 $+$ 关节速度 $n$ 维）。对基座（$\SE(3)$）用 Euler--Poincaré、对关节（$\mathbb T^n\cong\mathbb R^n$）用标准 Euler--Lagrange，合得 $(6+n)$ 维方程
\begin{equation}
  \mathbf M(\mathbf q)\dot{\mathbf v}+\mathbf C(\mathbf q,\mathbf v)\mathbf v+\mathbf g(\mathbf q)=\mathbf S^\top\boldsymbol\tau+\mathbf J_c^\top\boldsymbol\lambda,
  \label{eq:gm-floating-base}
\end{equation}
其中选择矩阵 $\mathbf S=[\mathbf 0_{n\times6}\;\;\mathbf I_n]$ 体现基座前 6 行不受驱（无电机直接作用于浮动基座），$\boldsymbol\tau\in\mathbb R^n$ 为关节力矩，$\mathbf J_c$ 为接触 Jacobian、$\boldsymbol\lambda$ 为接触力。质量矩阵分块
\begin{equation}
  \mathbf M(\mathbf q)=\begin{pmatrix}\mathbf M_{bb} & \mathbf M_{bj}\\ \mathbf M_{jb} & \mathbf M_{jj}\end{pmatrix},
  \label{eq:gm-M-block}
\end{equation}
$\mathbf M_{bb}\in\mathbb R^{6\times6}$ 为\emph{锁定惯性}（locked inertia，所有关节锁死时整机作为单刚体的 $6$D 惯性），$\mathbf M_{bj}=\mathbf M_{jb}^\top$ 为基座--关节耦合，$\mathbf M_{jj}$ 为关节空间质量矩阵。$\mathbf M_{bb}$ 依赖关节角 $\mathbf q_j$（改变关节角即改变各连杆相对基座的位置，从而改变整机绕基座的等效惯性）。

\begin{insight}
$\mathbf M_{bj}$ 耦合块解释了一个常见现象：某关节加速（$\ddot{\mathbf q}_j\ne\mathbf0$）时，通过 $\mathbf M_{bj}\ddot{\mathbf q}_j$ 在基座产生反力——这正是人形机器人挥臂会影响躯干平衡、空间机械臂关节运动会改变基座姿态（反作用力矩）的根源。\cref{eq:gm-floating-base} 的基座块 Coriolis $\mathbf C_{bb}$ 含两个来源：标准 Christoffel 项（来自 $\dot{\mathbf M}_{bb}$ 对 $\mathbf q_j$ 的依赖）与 Euler--Poincaré 的 $\operatorname{ad}^*_{\boldsymbol\xi_b}$ 项（体坐标系旋转效应）；在 \cref{ch:spatial-dynamics} 的 RNEA 中二者被自然合并，因递推已在体坐标系工作。
\end{insight}

\begin{remark}[$n_q\ne n_v$ 的深层原因与工程影响]
$\SO(3)$ 是 3 维流形，但无奇异表示（单位四元数）需 4 参数 $+1$ 约束，故位形维 $n_q$ 比速度维 $n_v$ 多 1（推广到 $\SE(3)$：$n_q=7+n$、$n_v=6+n$）。工程后果：(i) 不能用 $\dot{\mathbf q}$ 当切向量（须用 $\mathbf v\in\mathbb R^{n_v}$）；(ii) 各类 Jacobian 维度为 $?\times n_v$ 而非 $?\times n_q$；(iii) 积分须用李群指数（对基座部分），否则四元数模长偏离 1、旋转矩阵失正交；(iv) 位形差须返回切向量（$\log(\mathbf q_1^{-1}\mathbf q_2)\in\mathbb R^{n_v}$），两四元数的欧氏差无物理意义。接触力 $\mathbf J_c^\top\boldsymbol\lambda$ 作为额外广义力进入方程，约束力主体（KKT/Delassus/Baumgarte）留 \cref{ch:constrained-dynamics}；浮基动力学是 \cref{ch:reduced-models} 简化模型与腿足 MPC/WBC 的几何地基。
\end{remark}

% ════════════════════════════════════════════════════════════════
\section{质心动量、Centroidal 动力学与 Featherstone--Lie 群对照}
\label{sec:gm-centroidal}

\paragraph{质心动量矩阵与 Centroidal 动力学。}
浮基系统虽有数十自由度，其质心运动只有 $6$ 维。定义\emph{质心动量矩阵}（CMM）$\mathbf A_G(\mathbf q)\in\mathbb R^{6\times(6+n)}$，使质心动量
\begin{equation}
  \mathbf h=\begin{bmatrix}\mathbf k\\ \mathbf l\end{bmatrix}=\mathbf A_G(\mathbf q)\,\mathbf v\in\mathbb R^6,
  \label{eq:gm-cmm}
\end{equation}
其中 $\mathbf k$ 为关于质心的总角动量、$\mathbf l=m\mathbf v_{\mathrm{CoM}}$ 为总线动量。对其求时间导（Orin--Goswami--Lee 2013\cite{orin2013centroidal}）：
\begin{equation}
  \dot{\mathbf h}=\mathbf A_G(\mathbf q)\dot{\mathbf v}+\dot{\mathbf A}_G(\mathbf q,\mathbf v)\mathbf v=\sum_i\boldsymbol{\mathcal F}_{\mathrm{ext},i},
  \label{eq:gm-centroidal}
\end{equation}
即「线动量变化率 $=$ 合外力（含重力），角动量变化率 $=$ 合力矩」。本质上 Centroidal 动力学就是把整机等效为一个质点加陀螺的 Newton--Euler 方程，而 $\mathbf A_G$ 提供从关节空间精确控制质心动量的桥梁。

\begin{remark}[CMM 与关节空间惯量的关系及归属]
Orin \& Goswami（IROS 2008）引入了 CMM 的概念与结构；CMM 与关节空间惯量矩阵的关系（$\mathbf A_G$ 可由 $\mathbf M$ 前 6 行经与质心位置相关的变换抽取）及紧凑的 $O(n)$ 算法由 Orin, Goswami \& Lee（2013, \emph{Autonomous Robots} 35(2--3):161--176\cite{orin2013centroidal}）给出，意味着无需专门算法计算 CMM——RNEA/CRBA 的中间结果已含之。Wensing \& Orin（IJHR 2016）进一步给出工程改进的高效计算。Pinocchio 的 \texttt{ccrba}/\allowbreak\texttt{dccrba} 即此理论的工业实现，WBC/TSID 的角动量任务以 $\dot{\mathbf h}$ 为等式约束（主体留 \cref{ch:force-wbc}），本章只给 $\mathbf h,\dot{\mathbf h},\mathbf A_G$ 的几何起源。质心动量正是 $\SE(3)$ 动量映射在质心处的取值（\cref{sec:gm-momentum-map}），其守恒性将在 \cref{sec:gm-liepoisson} 经 Marsden--Weinstein 约化获得几何解释。
\end{remark}

\begin{table}[H]\centering\small
\caption{Featherstone 空间向量 $\leftrightarrow$ 李群几何精确对照（与 \cref{ch:spatial-dynamics}、\cref{ch:lie} 三向接线枢纽）}
\label{tab:gm-featherstone}
\begin{tabular}{lll}
\toprule
Featherstone 概念 & 李群几何对应 & 维度\\
\midrule
spatial velocity (motion) $\hat{\mathbf v}$ & body velocity $\boldsymbol\xi=\mathbf g^{-1}\dot{\mathbf g}\in\se(3)$ & 6\\
spatial force (wrench) $\hat{\mathbf f}$ & coadjoint momentum $\boldsymbol\mu\in\se(3)^*$ & 6\\
Plücker transform $\mathbf X$ & adjoint $\operatorname{Ad}_{\mathbf T}$ & $6\times6$\\
dual Plücker $\mathbf X^*=\mathbf X^{-\top}$ & coadjoint $\operatorname{Ad}^*_{\mathbf T^{-1}}$ & $6\times6$\\
motion cross $\times$ & adjoint action $\operatorname{ad}_{\boldsymbol\xi}$ & $6\times6$\\
force cross $\times^*$ & coadjoint action $-\operatorname{ad}^*_{\boldsymbol\xi}=-\operatorname{ad}_{\boldsymbol\xi}^\top$ & $6\times6$\\
spatial inertia $\boldsymbol{\mathcal I}$ & body inertia $\boldsymbol{\mathcal I}:\se(3)\to\se(3)^*$ & $6\times6$\\
RNEA 后向 $\hat{\mathbf f}=\boldsymbol{\mathcal I}\mathbf a+\hat{\mathbf v}\times^*(\boldsymbol{\mathcal I}\hat{\mathbf v})$ & EP on $\SE(3)$: $\boldsymbol{\mathcal I}\dot{\boldsymbol\xi}=\operatorname{ad}^*_{\boldsymbol\xi}(\boldsymbol{\mathcal I}\boldsymbol\xi)+\boldsymbol\tau$ & ---\\
\bottomrule
\end{tabular}
\end{table}

\begin{insight}
几何学家用抽象符号 $(\operatorname{ad},\operatorname{Ad},\operatorname{ad}^*)$，工程师用 $6\times6$ 矩阵 $(\times,\mathbf X,\times^*)$，几何内容毫无差别——如同用极坐标和笛卡尔坐标描述同一曲线。李群语言的优势在统一性与可推广性（从 $\SO(3)$ 到 $\SE(3)$ 到任意李群），Featherstone 记号的优势在计算效率与实现直觉。\cref{tab:gm-featherstone} 是本章与 \cref{ch:spatial-dynamics}、\cref{ch:lie} 三向接线的枢纽，与 J 章对照表保持完全一致。
\end{insight}

% ════════════════════════════════════════════════════════════════
% 第三部分 — 辛结构、约化与辛积分器
% ════════════════════════════════════════════════════════════════
\section{Hamilton 力学与辛流形}
\label{sec:gm-symplectic}

\paragraph{Hamilton 力学：从 $(\mathbf q,\dot{\mathbf q})$ 到 $(\mathbf q,\mathbf p)$。}
承 \cref{sec:gm-gravity}，对 $L=\tfrac12\dot{\mathbf q}^\top\mathbf M\dot{\mathbf q}-V$ 作 Legendre 变换：广义动量 $\mathbf p=\partial L/\partial\dot{\mathbf q}=\mathbf M(\mathbf q)\dot{\mathbf q}$。变换前提是 $L$ 对 $\dot{\mathbf q}$ 严格凸，等价于 $\mathbf M\succ0$（\cref{thm:gm-spd}），故恒可逆 $\dot{\mathbf q}=\mathbf M^{-1}\mathbf p$。Hamiltonian（总能量）
\begin{equation}
  H(\mathbf q,\mathbf p)=\mathbf p^\top\dot{\mathbf q}-L=\tfrac12\mathbf p^\top\mathbf M^{-1}(\mathbf q)\mathbf p+V(\mathbf q),
  \label{eq:gm-hamiltonian}
\end{equation}
正则方程 $\dot q_i=\partial H/\partial p_i$、$\dot p_i=-\partial H/\partial q_i+\tau_i$ 把 $n$ 个二阶 ODE 化为 $2n$ 个一阶 ODE。保守情形 $\dot H=\sum_i(\partial_{q_i}H\,\dot q_i+\partial_{p_i}H\,\dot p_i)=0$，机械能守恒。Poisson 括号 $\{F,G\}=\sum_i(\partial_{q_i}F\,\partial_{p_i}G-\partial_{p_i}F\,\partial_{q_i}G)$ 满足反对称、双线性、Leibniz、Jacobi 恒等式与基本关系 $\{q_i,p_j\}=\delta_{ij}$，且 $\dot F=\{F,H\}+\partial_t F$，故 $F$ 守恒 $\Leftrightarrow\{F,H\}=0$（不显含 $t$）。

\begin{definition}[辛流形]\label{def:gm-symplectic-manifold}
\emph{辛流形}是一对 $(M^{2n},\omega)$，$\omega$ 为 $M$ 上的 2-form 满足：(i) \emph{闭}，$\mathrm d\omega=0$；(ii) \emph{非退化}，$\omega^n\ne0$ 处处（等价于 $\omega^\flat:TM\to T^*M$ 同构）。余切丛 $T^*Q$ 天然辛：以 Liouville 1-form $\theta=p_i\,\mathrm dq^i$ 定义\emph{典则辛形式}
\begin{equation}
  \omega=-\mathrm d\theta=\sum_{i=1}^n\mathrm dq^i\wedge\mathrm dp_i.
  \label{eq:gm-canonical-form}
\end{equation}
\end{definition}

\begin{remark}[记号：辛标准矩阵 $\mathbf J$]
在 Darboux 坐标下 \cref{eq:gm-canonical-form} 对应辛标准矩阵 $\mathbf J=\bigl(\begin{smallmatrix}\mathbf0&\mathbf I\\-\mathbf I&\mathbf0\end{smallmatrix}\bigr)$。\textbf{此 $\mathbf J$ 为辛标准矩阵}，与 \cref{ch:lie} 的右雅可比 $\mathbf J_r$、运动学 Jacobian 同名而不同义；本章仅在辛节局部使用，动量映射另记 $\mathsf J$（\cref{sec:gm-momentum-map}）以区分。辛流形必为偶数维：反对称双线性形式在奇数维不可能非退化（$\det\mathbf A=\det(-\mathbf A^\top)=(-1)^n\det\mathbf A$，奇数 $n$ 得 $\det\mathbf A=0$）。闭性 $\mathrm d\omega=0$ 是 Hamilton 流保辛的全部物理（\cref{thm:gm-flow-symplectic}）。
\end{remark}

\begin{theorem}[Darboux]\label{thm:gm-darboux}
对任何辛流形 $(M,\omega)$ 与任一点 $x\in M$，存在局部坐标使 $\omega=\sum_i\mathrm dq^i\wedge\mathrm dp_i$。即\emph{辛几何无局部不变量}——所有辛流形局部都一样。
\end{theorem}

\begin{proof}[证明思路（Moser 路径法）]
令 $\omega_t=(1-t)\omega_0+t\omega$（$\omega_0=\sum\mathrm dq^i\wedge\mathrm dp_i$，选坐标使二者在 $x$ 点相等），小邻域内 $\omega_t$ 非退化。由闭性与 Poincaré 引理 $\omega-\omega_0=\mathrm d\alpha$。构造时变向量场 $X_t$ 使 $\iota_{X_t}\omega_t=-\alpha$，则 $\tfrac{\mathrm d}{\mathrm dt}\varphi_t^*\omega_t=\varphi_t^*(\mathcal L_{X_t}\omega_t+\dot\omega_t)=\varphi_t^*(\mathrm d\iota_{X_t}\omega_t+\omega-\omega_0)=\varphi_t^*(-\mathrm d\alpha+\mathrm d\alpha)=0$，故 $\varphi_1^*\omega=\omega_0$。
\end{proof}

\begin{insight}
Darboux 定理之于辛几何，如同「曲面局部可展为平面」之于欧氏几何——但黎曼几何有局部不变量（曲率），辛几何\emph{没有}。这解释了为何正则变换族极其丰富（不存在「辛曲率」限制坐标变换），也解释了为何辛几何的「有趣现象」（KAM 环面、Gromov 非挤压）都是\emph{整体}的。下表对照辛流形与黎曼流形：辛形式反对称、必偶数维、无局部不变量、保结构映射为辛同胚、刻画相空间结构；黎曼度量对称、任意维、有曲率张量、保结构映射为等距、刻画距离与角度。
\end{insight}

% ════════════════════════════════════════════════════════════════
\section{Hamilton 流保辛与正则变换}
\label{sec:gm-symplectic-flow}

\paragraph{Hamilton 方程 $=$ 辛梯度流。}
黎曼几何中梯度 $\nabla f$ 由度量定义（$g(\nabla f,X)=\mathrm df(X)$）；辛几何中「辛梯度」即 Hamilton 向量场 $X_H$ 由辛形式定义：
\begin{equation}
  \iota_{X_H}\omega=\mathrm dH.
  \label{eq:gm-hamiltonian-vf}
\end{equation}
非退化性保证 $X_H=(\omega^\flat)^{-1}(\mathrm dH)$ 唯一存在。在 Darboux 坐标下，设 $X_H=\dot q^i\partial_{q^i}+\dot p_i\partial_{p_i}$，代入 \cref{eq:gm-hamiltonian-vf} 得 $\dot q^i\,\mathrm dp_i-\dot p_i\,\mathrm dq^i=\partial_{q^i}H\,\mathrm dq^i+\partial_{p_i}H\,\mathrm dp_i$，比较系数即还原正则方程——它不是「规定」，而是辛结构的必然推论。

\begin{theorem}[Hamilton 流保辛]\label{thm:gm-flow-symplectic}
设 $\varphi_t$ 为 $X_H$ 的流，则 $\varphi_t^*\omega=\omega$ 对一切 $t$ 成立。
\end{theorem}

\begin{proof}
由 Cartan 魔术公式 $\mathcal L_X=\mathrm d\circ\iota_X+\iota_X\circ\mathrm d$：
\[
\mathcal L_{X_H}\omega=\mathrm d(\iota_{X_H}\omega)+\iota_{X_H}(\mathrm d\omega)=\mathrm d(\mathrm dH)+\iota_{X_H}(0)=0+0=0,
\]
第一项因 $\mathrm d^2=0$、第二项因 $\omega$ 闭。再由 $\tfrac{\mathrm d}{\mathrm dt}\varphi_t^*\omega=\varphi_t^*(\mathcal L_{X_H}\omega)=0$ 得结论。
\end{proof}

\begin{insight}
这一行证明是整个经典力学的几何灵魂：能量守恒、Liouville 体积守恒、Poincaré 回归、不存在 Hamilton 吸引子——全是 $\mathcal L_{X_H}\omega=0$ 的推论，而证明只用了两个事实 $\mathrm d^2=0$ 与 $\mathrm d\omega=0$。其直接推论是 \emph{Liouville 定理}：$\varphi_t^*(\omega^n/n!)=(\varphi_t^*\omega)^n/n!=\omega^n/n!$，相空间体积守恒（「相空间不可压缩」），故耗散系统（有吸引子）必非 Hamilton。\textbf{需厘清}：保辛 $\ne$ 保面积——4 维相空间中 $\omega=\mathrm dq^1\wedge\mathrm dp_1+\mathrm dq^2\wedge\mathrm dp_2$ 可在某 2 维子空间取零，真正的不变量是 $\omega$（一个 2-form）本身；保体积是保辛的推论而非定义。
\end{insight}

\paragraph{正则变换。}
微分同胚 $\varphi$ 称\emph{辛同胚}（正则变换）若 $\varphi^*\omega=\omega$，在 Darboux 坐标下等价于 Jacobian 满足辛矩阵条件 $\mathbf M^\top\mathbf J\mathbf M=\mathbf J$（即辛群 $\mathrm{Sp}(2n,\mathbb R)$ 的定义）\footnote{史料常把 $\mathbf M^\top\mathbf J\mathbf M=\mathbf J$ 这一现代记号命名为「辛矩阵条件」并归于 1838 年 Jacobi。「辛」（symplectic）一词实由 Weyl 1939 年在 \emph{The Classical Groups} 创造（源希腊语「编织在一起」，是 complex 的 calque），以替代旧名 \emph{line complex group}；中文「辛」为华罗庚所译。把现代命名前置到 1838 属时代错置——Jacobi 在 1837 前后系统化的是正则变换理论（其变换保持正则方程形式，即后世所谓的辛条件），技术内容无误。}。正则变换可经四类生成函数 $F_1(\mathbf q,\mathbf Q),F_2(\mathbf q,\mathbf P),F_3(\mathbf p,\mathbf Q),F_4(\mathbf p,\mathbf P)$ 构造——每个辛积分器都对应一个生成函数（\cref{sec:gm-integrators}）。Poincaré 积分不变量 $\oint_\gamma\mathbf p\cdot\mathrm d\mathbf q$ 沿 Hamilton 流的闭曲线像不变，是 action-angle 变量与 WKB 近似的基础。

% ════════════════════════════════════════════════════════════════
\section{动量映射与 Noether 定理的辛版本}
\label{sec:gm-momentum-map}

\paragraph{动机：守恒量的几何本质。}
\cref{thm:gm-el} 旁的 Noether 定理（Lagrange 版）给出守恒量，却不揭示守恒量之间的关系（如角动量分量的 Poisson 括号）。辛版本经\emph{动量映射}统一一切。设李群 $G$ 辛作用于 $(P,\omega)$，每个 $\boldsymbol\xi\in\mathfrak g$ 生成无穷小作用 $\boldsymbol\xi_P(p)=\tfrac{\mathrm d}{\mathrm dt}|_0\Phi_{\exp(t\boldsymbol\xi)}(p)$。

\begin{definition}[动量映射]\label{def:gm-momentum-map}
光滑映射 $\mathsf J:P\to\mathfrak g^*$ 称作该作用的\emph{动量映射}，若对每个 $\boldsymbol\xi\in\mathfrak g$，分量函数 $\mathsf J_{\boldsymbol\xi}(p):=\langle\mathsf J(p),\boldsymbol\xi\rangle$ 满足
\begin{equation}
  \iota_{\boldsymbol\xi_P}\omega=\mathrm d\mathsf J_{\boldsymbol\xi},
  \label{eq:gm-momentum-map}
\end{equation}
即 $\boldsymbol\xi_P$ 是由 $\mathsf J_{\boldsymbol\xi}$ 生成的 Hamilton 向量场。
\end{definition}

标准例：(i) 平移群 $G=\mathbb R^3$ 作用于 $T^*\mathbb R^3$，$\iota_{\boldsymbol\xi_P}(\mathrm dq^i\wedge\mathrm dp_i)=\xi^i\mathrm dp_i=\mathrm d\langle\mathbf p,\boldsymbol\xi\rangle$，故 $\mathsf J=\mathbf p$（线动量）；(ii) $G=\SO(3)$ 作用 $\Phi_{\mathbf R}(\mathbf q,\mathbf p)=(\mathbf R\mathbf q,\mathbf R\mathbf p)$，得 $\mathsf J=\mathbf q\times\mathbf p$（角动量）；(iii) $G=\SE(3)$ 作用于 $T^*\SE(3)$，$\mathsf J=$（线动量，关于原点的角动量）$\in\mathbb R^6$——这正是 \cref{sec:gm-centroidal} 质心动量的几何起源。

\begin{theorem}[Noether 定理辛版本]\label{thm:gm-noether-symplectic}
若 Hamiltonian $H$ 在 $G$ 作用下不变（$H\circ\Phi_g=H$），则动量映射 $\mathsf J$ 沿 $X_H$ 的流守恒：$\tfrac{\mathrm d}{\mathrm dt}\mathsf J(\varphi_t(p))=0$。
\end{theorem}

\begin{proof}
$\{\mathsf J_{\boldsymbol\xi},H\}=\omega(X_{\mathsf J_{\boldsymbol\xi}},X_H)=\omega(\boldsymbol\xi_P,X_H)=(\iota_{\boldsymbol\xi_P}\omega)(X_H)=\mathrm d\mathsf J_{\boldsymbol\xi}(X_H)=\mathcal L_{X_H}\mathsf J_{\boldsymbol\xi}$。另一方面 $\mathcal L_{X_H}\mathsf J_{\boldsymbol\xi}=\tfrac{\mathrm d}{\mathrm dt}\mathsf J_{\boldsymbol\xi}(\varphi_t)$，而由 $G$-不变性 $\mathcal L_{\boldsymbol\xi_P}H=\mathrm dH(\boldsymbol\xi_P)=0$。故 $\dot{\mathsf J}_{\boldsymbol\xi}=0$，$\boldsymbol\xi$ 任意得 $\dot{\mathsf J}=\mathbf0$。
\end{proof}

\paragraph{存在性与等变性。}
不是所有辛作用都有动量映射：充分条件含 $H^1(M,\mathbb R)=0$（每个闭 1-form 恰当）、$G$ 半单（Whitehead 引理）、或\textbf{ $T^*Q$ 的自然提升作用——它总有等变动量映射}。机器人相空间 $T^*Q$ 是余切丛、群作用来自位形空间的自然作用，故动量映射总存在且等变。\emph{等变性} $\mathsf J(\Phi_g\cdot p)=\operatorname{Ad}^*_{g^{-1}}\mathsf J(p)$ 保证守恒量的 Poisson 括号与李代数结构常数一致：$\{\mathsf J_{\boldsymbol\xi},\mathsf J_{\boldsymbol\eta}\}=\mathsf J_{[\boldsymbol\xi,\boldsymbol\eta]}$。

\begin{insight}
角动量分量满足 $\{L_1,L_2\}=L_3$（及轮换）——这正是 $\so(3)$ 的结构常数！动量映射的等变性就是这一关系的坐标无关版本：Poisson 括号把相空间的代数结构与对称群的李代数结构双向绑定。其机器人含义见对称破缺表（\cref{tab:gm-noether}）：自由刚体（无重力）$\SE(3)$ 不变 $\Rightarrow$ 6 个空间动量守恒；有竖直重力时退化到 $\SE(2)$（绕竖轴旋转 $+$ 两个水平平移，共 3 维）$\Rightarrow$ 只剩 z 角动量 $+$ 水平线动量（3 个，竖直线动量\emph{不}守恒，因重力沿 z 做功）；加平面约束退到 $\SO(2)$ $\Rightarrow$ 仅 z 角动量。守恒量不是偶然，而是对称性的必然——这帮助控制器设计预判哪些量无需主动控制（如飞行相的角动量守恒可用于腿足跳跃/翻滚规划）。
\end{insight}

\begin{table}[H]\centering\small
\caption{对称破缺与守恒量丧失}
\label{tab:gm-noether}
\begin{tabular}{lll}
\toprule
对称性 & 不变群 & 守恒量（维数）\\
\midrule
完全自由（无重力）& $\SE(3)$ & 6 个空间动量（6）\\
有竖直重力 & $\SE(2)$ & z 角动量 $+$ 水平线动量（3）\\
有重力 $+$ 平面约束 & $\SO(2)$ & z 角动量（1）\\
\bottomrule
\end{tabular}
\end{table}

% ════════════════════════════════════════════════════════════════
\section{Lie--Poisson 约化、Casimir 与 Marsden--Weinstein 定理}
\label{sec:gm-liepoisson}

\paragraph{Poisson 流形与 Lie--Poisson 结构。}
\emph{Poisson 流形} $(P,\{\cdot,\cdot\})$ 比辛流形更一般（无需偶数维、无需非退化），括号满足反对称、Leibniz、Jacobi。辛流形皆 Poisson（$\{F,G\}=\omega(X_F,X_G)$），反之不然。李代数对偶 $\mathfrak g^*$ 上有典则\emph{Lie--Poisson 括号}
\begin{equation}
  \{F,G\}_\pm(\boldsymbol\mu)=\pm\Big\langle\boldsymbol\mu,\Big[\frac{\delta F}{\delta\boldsymbol\mu},\frac{\delta G}{\delta\boldsymbol\mu}\Big]\Big\rangle,
  \label{eq:gm-lie-poisson}
\end{equation}
$+/-$ 对应右/左不变约化（机器人多用左约化，取「$-$」）。对任意 $H:\mathfrak g^*\to\mathbb R$，Lie--Poisson 方程 $\dot{\boldsymbol\mu}=\operatorname{ad}^*_{\delta H/\delta\boldsymbol\mu}\boldsymbol\mu$，经 Legendre 变换 $H(\boldsymbol\mu)=\langle\boldsymbol\mu,\boldsymbol\xi\rangle-\ell(\boldsymbol\xi)$ 与 Euler--Poincaré 等价。$\SO(3)$ 上取 $\boldsymbol\mu=\boldsymbol\pi$、$H=\tfrac12\boldsymbol\pi^\top\mathbf I_c^{-1}\boldsymbol\pi$、$\delta H/\delta\boldsymbol\pi=\boldsymbol\omega$，得 $\dot{\boldsymbol\pi}=\boldsymbol\pi\times\boldsymbol\omega$——又是 Euler 方程，但现在表达在纯动量空间 $\mathbb R^3$，无需显式位形 $\mathbf R$。

\begin{definition}[Casimir 函数]\label{def:gm-casimir}
$C\in C^\infty(P)$ 称\emph{Casimir}，若 $\{C,F\}=0$ 对一切 $F$。它对\emph{任何} Hamiltonian 都守恒——是 Poisson 结构本身的约束，不依赖具体 $H$。
\end{definition}

$\SO(3)$ 上 $C(\boldsymbol\pi)=|\boldsymbol\pi|^2$ 是 Casimir（$\{C,F\}_-=-\boldsymbol\pi\cdot(2\boldsymbol\pi\times\nabla F)=-2(\boldsymbol\pi\times\boldsymbol\pi)\cdot\nabla F=0$）：无论惯量张量如何，体角动量模长总守恒。$\SE(3)^*$ 有两个 Casimir：$C_1=\mathbf m\cdot\mathbf f$（螺旋不变量）与 $C_2=|\mathbf f|^2$（线动量模平方），故其余伴随轨道一般为 4 维（$6-2$）。Kirchhoff 方程——自由刚体在理想流体中的运动——正是 $\se(3)^*$ 上的 Lie--Poisson 方程 $\dot{\mathbf f}=\mathbf f\times\boldsymbol\omega,\,\dot{\boldsymbol\pi}=\boldsymbol\pi\times\boldsymbol\omega+\mathbf f\times\mathbf v$，是 $\SE(3)$ Lie--Poisson 结构的物理实例。

\begin{theorem}[Kirillov--Kostant--Souriau：辛叶 $=$ 余伴随轨道]\label{thm:gm-kks}
Poisson 流形 $\mathfrak g^*$ 可分层为\emph{辛叶}，每层上 Poisson 括号由某辛结构诱导，且辛叶恰是\emph{余伴随轨道} $\mathcal O_{\boldsymbol\mu}=\{\operatorname{Ad}^*_g\boldsymbol\mu:g\in G\}$。$\SO(3)^*\cong\mathbb R^3$ 的辛叶是各半径同心球面 $S^2_r$（$\operatorname{Ad}^*$ 即标准旋转作用）。
\end{theorem}

\begin{insight}[Poinsot 构造与中间轴不稳定]
Casimir $|\boldsymbol\pi|^2$ 定义辛叶（球面 $S^2_r$），Hamiltonian $H=\tfrac12\boldsymbol\pi^\top\mathbf I_c^{-1}\boldsymbol\pi$ 的等值面（能量椭球）与球面的交线即 Euler 刚体的轨迹——这正是 Poinsot 构造的现代几何解释（2D 辛叶上的 1D 等能面即轨迹）。由此读出经典结论：惯量 $I_1>I_2>I_3$ 时，绕最大轴 $I_1$、最小轴 $I_3$ 旋转稳定，\textbf{绕中间轴 $I_2$ 旋转不稳定}——能量椭球与动量球面在中间轴附近交出的是\emph{鞍点}（异宿轨道连接两个不稳定平衡），轨迹半绕球面后翻转，即网球拍定理／Dzhanibekov 效应（Poinsot 1834 首述，Dzhanibekov 1985 太空发现）。线性化验证：对平衡 $\boldsymbol\omega=(0,\Omega,0)$ 取小扰动得 $A C\ddot\epsilon_1=(B-C)(A-B)\Omega^2\epsilon_1$，$A<B<C$ 时系数为正，指数增长。Explorer 1（1958）因柔性天线耗能、在 $|\boldsymbol\pi|^2$ 守恒下趋向最小动能态，从最小惯量轴自发翻到最大惯量轴，是此现象的工程实例。
\end{insight}

\begin{theorem}[$T^*G/G\cong\mathfrak g^*$ 与 Marsden--Weinstein 辛约化]\label{thm:gm-reduction}
作为 Poisson 流形，$T^*G/G\cong\mathfrak g^*$（Marsden--Ratiu Thm 13.1.1）：把 $T^*G$（辛，$2\dim G$ 维）关于 $G$ 左作用取商得 $\mathfrak g^*$（Poisson，$\dim G$ 维）。更一般地，设 $G$ Hamilton 等变作用于辛 $(P,\omega)$、动量映射 $\mathsf J$、$\boldsymbol\mu$ 为正则值、$G_{\boldsymbol\mu}$ 为其迷向子群且自由真作用于 $\mathsf J^{-1}(\boldsymbol\mu)$，则
\begin{equation}
  P_{\boldsymbol\mu}:=\mathsf J^{-1}(\boldsymbol\mu)/G_{\boldsymbol\mu}
  \label{eq:gm-marsden-weinstein}
\end{equation}
是辛流形，维度 $\dim P_{\boldsymbol\mu}=\dim P-\dim G-\dim G_{\boldsymbol\mu}$（Marsden--Weinstein 1974\cite{marsden1974reduction}）。
\end{theorem}

\begin{insight}
对 $\SO(3)$ 刚体：完整相空间 $T^*\SO(3)$（6 维），左 $\SO(3)$ 作用（改变空间坐标系不影响体坐标物理）约化到 $\so(3)^*\cong\mathbb R^3$（3 维），约化方程即 Euler 方程——空间坐标信息 $\mathbf R$ 被「商掉」，因自由刚体运动不依赖其在空间中的朝向。对浮基人形 $Q=\SE(3)\times\mathbb R^n$、$G=\SE(3)$，无外力时 $\boldsymbol\mu=\mathbf0$，$\dim P_0=2(6+n)-6-6=2n$——约化后只剩 $n$ 个内部自由度，\textbf{这就是「浮基消失」的几何原因}，而 Centroidal 动力学（\cref{sec:gm-centroidal}）正是这一约化的局部坐标表达。当重力破坏 $\SE(3)$ 完全不变性时，引入 advected parameter $\boldsymbol\Gamma=\mathbf R^\top\mathbf g/|\mathbf g|$（体坐标系重力方向），得半直积 Lie--Poisson 的 heavy top 方程 $\dot{\boldsymbol\pi}=\boldsymbol\pi\times\boldsymbol\omega+mgl\,\boldsymbol\Gamma\times\mathbf e_3,\,\dot{\boldsymbol\Gamma}=\boldsymbol\Gamma\times\boldsymbol\omega$（Marsden--Ratiu--Weinstein 1984\cite{marsden1984semidirect}）。
\end{insight}

% ════════════════════════════════════════════════════════════════
\section{辛积分器与修正 Hamiltonian}
\label{sec:gm-integrators}

\paragraph{动机：长时不漂胜过单步更准。}
用 RK4 积分 Kepler 问题，$10^6$ 步后轨道系统性收缩（行星螺旋坠入太阳）；用 Störmer--Verlet，$10^9$ 步后仍在正确能量面上振荡。原因不是精度（RK4 是 4 阶、Verlet 仅 2 阶），而是\emph{保辛性}。一步法 $\Phi_h$ 称\emph{辛的}，若 $\Phi_h^*\omega=\omega$，等价于 Jacobian 满足 $(\partial\Phi_h/\partial\mathbf y)^\top\mathbf J(\partial\Phi_h/\partial\mathbf y)=\mathbf J$。

\begin{theorem}[Symplectic Euler 与 Störmer--Verlet 保辛]\label{thm:gm-symplectic-euler}
对可分 Hamiltonian $H=T(\mathbf p)+V(\mathbf q)$，\emph{辛 Euler} $\mathbf p_{n+1}=\mathbf p_n-h\nabla V(\mathbf q_n),\,\mathbf q_{n+1}=\mathbf q_n+h\nabla T(\mathbf p_{n+1})$（先更新 $\mathbf p$，再用 $\mathbf p_{n+1}$ 更新 $\mathbf q$）保辛；\emph{Störmer--Verlet}（半步—全步—半步）作为两个辛 Euler 的复合亦保辛，且时间可逆、二阶、对可分 $H$ 显式。
\end{theorem}

\begin{proof}[证明骨架]
辛 Euler 的 Jacobian $\mathbf M=\bigl(\begin{smallmatrix}\mathbf I-h^2 T_{pp}V_{qq} & h T_{pp}\\ -h V_{qq} & \mathbf I\end{smallmatrix}\bigr)$（取相空间序 $(\mathbf q,\mathbf p)$：$\mathbf q_{n+1}=\mathbf q_n+h\nabla T(\mathbf p_n-h\nabla V(\mathbf q_n))$ 对 $\mathbf q_n$ 求导即得左上块 $\mathbf I-h^2T_{pp}V_{qq}$），直接代入并用 $V_{qq},T_{pp}$ 对称验证 $\mathbf M^\top\mathbf J\mathbf M=\mathbf J$（左上块取负号方使该恒等式成立，正号则 $\mathbf M^\top\mathbf J\mathbf M-\mathbf J$ 不为零）；亦可由 $F_1$ 型生成函数 $S(\mathbf q_n,\mathbf q_{n+1})$ 与辛变换一一对应得辛性。Störmer--Verlet 是辛 Euler（步长 $h/2$）与其伴随的复合，两个辛映射的复合仍辛。
\end{proof}

\begin{theorem}[向后误差分析：修正 Hamiltonian]\label{thm:gm-backward-error}
对辛一步法 $\Phi_h$ 存在形式级数 \emph{modified Hamiltonian} $\tilde H(\mathbf y;h)=H+h^r H_{r+1}+h^{r+1}H_{r+2}+\cdots$，使 $\Phi_h$ 在形式级数意义下\emph{精确求解} $\dot{\mathbf y}=X_{\tilde H}(\mathbf y)$（Hairer--Lubich--Wanner\cite{hairer2006geometric}）。对解析 $H$ 与 $h<h^*$，能量误差 $|H(\mathbf y_n)-H(\mathbf y_0)|\le C\mathrm e^{-c/h}+C h^r t_n$，直到指数长时 $t\le\mathrm e^{c/h}$。
\end{theorem}

\begin{insight}
这解释了「辛保辛却不保能量」的真因：辛积分器精确保的是\emph{辛结构}，对原始 $H$ 的能量误差是\emph{有界振荡}（$O(h^r)$）而非\emph{单调漂移}——因为它精确求解的是修正 Hamiltonian $\tilde H=H+O(h^r)$。如同一只「略有偏差却齿轮比精确」的钟，误差是有界周期振荡而非累积。正因此 Verlet 可做 $10^9$ 步行星仿真而不失真。\textbf{需厘清}：不可能同时精确保辛与精确保能量（除非精确解，Ge--Marsden 1988）；若须精确保能量则用 discrete gradient 方法，但它不保辛。
\end{insight}

\paragraph{其他辛方法与 Lie 群变分积分器。}
\emph{隐式中点法} $(\mathbf y_{n+1}-\mathbf y_n)/h=X_H((\mathbf y_n+\mathbf y_{n+1})/2)$ 是唯一对\emph{任意} $H$ 都辛的 1 级 Gauss--Legendre 方法——当 Coriolis 使 $H=\tfrac12\mathbf v^\top\mathbf M(\mathbf q)\mathbf v+V$ 不可分时，它是唯一可用的辛 RK（代价是隐式）。\emph{Yoshida 4 阶}由三次 Verlet 以步长 $(w_1,w_2,w_1)$ 复合，$w_1=1/(2-2^{1/3}),\,w_2=1-2w_1<0$（负步长「时间倒退一小段」是达到高阶的代价）\cite{yoshida1990construction}。\emph{Lie 群变分积分器}（Marsden--West 2001\cite{marsden2001discrete}）直接在李群上写离散 Hamilton 原理：对 $\SO(3)$ 刚体取 $\mathbf R_{k+1}=\mathbf R_k\mathbf F_k$（$\mathbf F_k\in\SO(3)$），自动保 $\mathbf R_k\in\SO(3)$（无须重投影）、保辛、由离散 Noether 定理\emph{精确}保角动量、长时能量漂移 $O(h^2)$ 有界——优于「普通 RK $+$ 单位化」（后者破坏动量守恒与辛性）。

\begin{remark}[归属与库的辛性]
辛差分格式由冯康（Feng Kang）1985 年在 1984 北京微分几何与微分方程国际会议上独立提出\cite{feng1985difference}，系统建立了「保结构算法」理念与中国辛算法学派（专著见冯康、秦孟兆\cite{feng2003symplectic}），几乎与 Ruth（1983，正则积分技术，3 阶\cite{ruth1983canonical}）同期；4 阶高阶辛分裂为 Forest--Ruth（1990\cite{forest1990fourth}）与 Yoshida（1990）。主流多体库（MuJoCo/Pinocchio/Drake）默认积分器因 Coriolis 与接触不严格保辛，仅作该理论的工业近似。约束系统的辛积分器（SHAKE/RATTLE）见 \cref{ch:constrained-dynamics}。
\end{remark}

% ════════════════════════════════════════════════════════════════
\section{辛几何与最优控制：Pontryagin 辛表述}
\label{sec:gm-pontryagin}

\paragraph{最优控制 $=$ 余切丛上的 Hamilton 流。}
\cref{ch:opt-variational-pmp} 给出 Pontryagin 极小值原理的算法主体，本节给其几何观点。对问题 $\min_{\mathbf u}\int_0^T L(\mathbf x,\mathbf u)\,\mathrm dt$ s.t. $\dot{\mathbf x}=\mathbf f(\mathbf x,\mathbf u)$，定义 Pontryagin Hamiltonian
\begin{equation}
  H(\mathbf x,\boldsymbol\lambda,\mathbf u)=\langle\boldsymbol\lambda,\mathbf f(\mathbf x,\mathbf u)\rangle-L(\mathbf x,\mathbf u),
  \label{eq:gm-pontryagin-h}
\end{equation}
$\boldsymbol\lambda\in T^*_{\mathbf x}M$ 为共态。状态方程 $\dot{\mathbf x}=\partial H/\partial\boldsymbol\lambda$ 与协态方程 $\dot{\boldsymbol\lambda}=-\partial H/\partial\mathbf x$ 合起来即 $T^*M$ 上的 Hamilton 正则方程：代入最优 $\mathbf u^*(\mathbf x,\boldsymbol\lambda)$ 后，最优轨道 $(\mathbf x^*(t),\boldsymbol\lambda^*(t))$ 是约化 Hamiltonian $H^*$ 的 Hamilton 流。

\begin{insight}
共态 $\boldsymbol\lambda$ 不是数值技巧，而是余切向量（「动量」）：$\boldsymbol\lambda(t)=\partial V/\partial\mathbf x$（$V$ 为值函数），衡量「状态偏离一点、最优代价变化多少」。由此 DDP/iLQR 的 backward pass 中 Riccati 方程是\emph{辛流的线性化}（线性 Hamilton 系统有 Riccati 解），$V_{\mathbf{xx}}$ 是 $T^*M$ 上 Lagrangian 子流形（最优回报流形）的曲率——这是 \cref{ch:ddp-ilqr} 几何性质的根源。DMOC（离散力学与最优控制，Kobilarov--Marsden 2011\cite{kobilarov2011discrete}）把变分积分器嵌入最优控制得「离散 PMP」：离散 Noether 自动保守恒量、离散辛结构使非线性规划 Hessian 结构更好。\cref{ch:opt-variational-pmp}（变分/PMP）与 \cref{ch:ddp-ilqr}（DDP/iLQR）讲算法主体，本章提供「最优控制 $=$ 辛流」的几何观点；$\SE(3)$ 上的几何跟踪控制见 Lee--Leok--McClamroch 2010\cite{lee2010geometric}。
\end{insight}

% ════════════════════════════════════════════════════════════════
\section{进阶专题、学习与采样中的回响，及全章脉络}
\label{sec:gm-frontier}

\paragraph{辛结构在学习与采样中的回响。}
本章的辛思想在机器学习中有三个概念锚点（公式推导与库实现留学习部）。其一，\emph{Hamiltonian Monte Carlo}（HMC）以 Störmer--Verlet 在 $T^*\Theta$ 上积分：因 Hamilton 流保相空间体积（Liouville 定理），保证 detailed balance；用非辛积分器会系统偏倚采样分布，Stan/NumPyro 的 NUTS 即基于辛积分。其二，\emph{自然梯度}（Amari 1998\cite{amari1998natural}）$\dot{\boldsymbol\theta}=-\mathbf F^{-1}\nabla L$ 是 $T^*\Theta$ 上 Hamiltonian $H=L+\tfrac12\mathbf p^\top\mathbf F^{-1}\mathbf p$（Fisher 度量定义动能）在强阻尼极限下的投影，Natural Policy Gradient（Kakade 2001\cite{kakade2001natural}）与 TRPO 的 Fisher 信任域由此而来。其三，\emph{辛网络}把辛结构作为归纳偏置：Hamiltonian Neural Networks（Greydanus, Dzamba \& Yosinski 2019\cite{greydanus2019hamiltonian}）学一个 $H_\theta(\mathbf q,\mathbf p)$ 再以辛积分器积分，SympNets（Jin, Zhang, Zhu, Tang \& Karniadakis 2020\cite{jin2020sympnets}）堆叠辛层使整体辛——二者均使长 rollout 不漂移能量。Fisher 信息与 SLAM 因子图 Hessian 的联系一句接 \cref{ch:isam2}（与 \cref{ch:spatial-dynamics} 的「ABA $\leftrightarrow$ $\mathbf L\mathbf D\mathbf L^\top$ $\leftrightarrow$ 消元树」同为跨部亮点：动力学/辛的几何结构即估计后端的几何结构）。

\paragraph{进阶纯数学指针。}
若干进阶专题各有真实机器人接口（范围所限只给定性指针）：(i) Routh reduction（循环坐标消去）$\leftrightarrow$ 带反作用轮卫星的本体姿态 $+$ 轮角动量降维；(ii) Arnold--Liouville 可积性与 KAM 定理 $\leftrightarrow$ 周期步态为退化 KAM 环面、稳定域为环面残存、Arnold diffusion 对应「摔倒」；(iii) Poincaré 回归与回归映射 $\leftrightarrow$ 步态 return map——辛 Poincaré 映射的特征值必成对 $(\lambda,1/\lambda)$，故\emph{不存在渐近稳定的被动步态}（被动行走只能边际稳定，渐近稳定须加耗散或控制，二者皆破坏辛结构）；(iv) Gromov 非挤压定理 $\leftrightarrow$ Capture Region 的辛容量上界（辛积分器保辛容量，不会虚假地扩大可达区域）；(v) Dirac 结构（Yoshimura--Marsden 2006）$\leftrightarrow$ 接触瞬间破坏连续辛结构时的统一框架。零角动量约束下「内部形状变化 $\to$ 外部位移」的纤维丛 holonomy（几何相位）则连接 locomotion 前沿（蛇形/猫翻身/卫星姿态重定向）。

\begin{insight}[本章脉络全景]
本章把 $\mathbf M\ddot{\mathbf q}+\mathbf C\dot{\mathbf q}+\mathbf g=\boldsymbol\tau$ 沿三条主线走通：分析力学（变分原理 $\to$ $\mathbf M,\mathbf C,\mathbf g$ $\to$ $\dot{\mathbf M}-2\mathbf C$ 反对称 $\to$ 被动性）、几何力学（左不变约化 $\to$ Euler--Poincaré $\to$ 浮基 $\to$ Centroidal）、辛结构（辛流形 $\to$ 保辛 $\to$ 动量映射/Noether $\to$ Lie--Poisson/Marsden--Weinstein $\to$ 辛积分器）。贯穿三部分的主梁是同一个 Coriolis：Christoffel $c_{ijk}$（\cref{def:gm-christoffel}）$=$ 力叉积 $\mathbf v\times^*(\boldsymbol{\mathcal I}\mathbf v)$（\cref{ch:spatial-dynamics}）$=$ $\operatorname{ad}^*_{\boldsymbol\xi}\boldsymbol\mu$（\cref{thm:gm-euler-poincare}）$=$ Lie--Poisson 的 $\operatorname{ad}^*$（\cref{sec:gm-liepoisson}）。向后，本章是 \cref{ch:constrained-dynamics}（约束/解析微分）、\cref{ch:reduced-models}（简化模型）、\cref{ch:contact-mechanics}（接触）、\cref{ch:operational-space}（操作空间）、\cref{ch:force-wbc}（WBC）的几何地基，也为机械臂部（固定基串联臂的 $\mathbf M,\mathbf C,\mathbf g$ 特化）、腿足部（浮基/Centroidal/中间轴/飞行相动量守恒）、估计部（\cref{ch:isam2} 的 $SE_2(3)$/InEKF group-affine 一致性\cite{barrau2017invariant}\cite{hartley2020contact}）提供分析力学根基。
\end{insight}

% ════════════════════════════════════════════════════════════════
\section*{本章小结}
\addcontentsline{toc}{section}{本章小结}

本章把机器人标准运动方程 $\mathbf M\ddot{\mathbf q}+\mathbf C\dot{\mathbf q}+\mathbf g=\boldsymbol\tau$ 从「一个公式」升级为三条殊途同归的几何叙事。\textbf{分析力学}（\crefrange{sec:gm-history}{sec:gm-gravity}）：由 d'Alembert 原理与 Hamilton 原理推出 Euler--Lagrange 方程，构造对称正定的 $\mathbf M$（\cref{thm:gm-spd}）、Christoffel 符号与作为位形空间 Levi--Civita 联络的 Coriolis 矩阵（\cref{thm:gm-levicivita}），证明 $\dot{\mathbf M}-2\mathbf C$ 反对称的强/弱版本（\cref{thm:gm-skew}）及其即时推论被动性（\cref{thm:gm-passivity}）。\textbf{几何力学}（\crefrange{sec:gm-why-geom}{sec:gm-centroidal}）：以左不变约化与约束变分 $\delta\boldsymbol\xi=\dot{\boldsymbol\eta}+\operatorname{ad}_{\boldsymbol\xi}\boldsymbol\eta$ 推出 Euler--Poincaré 方程（\cref{thm:gm-euler-poincare}），退化得 Euler 刚体方程与 $6$D Newton--Euler，扩展到浮基 $\mathbf M\dot{\mathbf v}+\mathbf C\mathbf v+\mathbf g=\mathbf S^\top\boldsymbol\tau+\mathbf J_c^\top\boldsymbol\lambda$ 与 Centroidal 动力学。\textbf{辛结构}（\crefrange{sec:gm-symplectic}{sec:gm-frontier}）：辛流形与 Darboux 定理、Hamilton 流保辛的一行证明（\cref{thm:gm-flow-symplectic}）、动量映射与 Noether 辛版（\cref{thm:gm-noether-symplectic}）、Lie--Poisson 约化与 Marsden--Weinstein 定理（\cref{thm:gm-reduction}）、辛积分器与修正 Hamiltonian（\cref{thm:gm-backward-error}）。

\begin{table}[H]\centering\small
\caption{本章核心结论速查}
\label{tab:gm-summary}
\begin{tabular}{lll}
\toprule
结论 & 一句话 & 对应\\
\midrule
Euler--Lagrange & $\delta S=0\Rightarrow\tfrac{\mathrm d}{\mathrm dt}\partial_{\dot q}L-\partial_q L=\tau$ & \cref{thm:gm-el}\\
$\mathbf M\succ0$ & 无零能量运动、$\mathbf M^{-1}$ 存在 & \cref{thm:gm-spd}\\
Christoffel $=$ Levi--Civita & 无外力运动 $=$ 测地线 & \cref{thm:gm-levicivita}\\
$\dot{\mathbf M}-2\mathbf C$ 反对称 & 强版仅 Christoffel；弱版任意 $\mathbf C$ & \cref{thm:gm-skew}\\
被动性 & $\dot H=\dot{\mathbf q}^\top\boldsymbol\tau$ & \cref{thm:gm-passivity}\\
Euler--Poincaré & $\tfrac{\mathrm d}{\mathrm dt}\partial_{\boldsymbol\xi}\ell=\operatorname{ad}^*_{\boldsymbol\xi}\partial_{\boldsymbol\xi}\ell+\boldsymbol\tau$ & \cref{thm:gm-euler-poincare}\\
SE(3) 退化 & $6$D Newton--Euler $=$ EP on $\SE(3)$ & \cref{eq:gm-se3-newton-euler}\\
Hamilton 流保辛 & $\mathcal L_{X_H}\omega=\mathrm d(\mathrm dH)+\iota_{X_H}\mathrm d\omega=0$ & \cref{thm:gm-flow-symplectic}\\
Noether 辛版 & $G$-不变 $H\Rightarrow\dot{\mathsf J}=0$ & \cref{thm:gm-noether-symplectic}\\
Lie--Poisson & $\dot{\boldsymbol\mu}=\operatorname{ad}^*_{\delta H/\delta\boldsymbol\mu}\boldsymbol\mu$；Casimir $|\boldsymbol\pi|^2$ & \cref{thm:gm-kks}\\
Marsden--Weinstein & $P_{\boldsymbol\mu}=\mathsf J^{-1}(\boldsymbol\mu)/G_{\boldsymbol\mu}$；浮基 $\mu{=}0\Rightarrow2n$ & \cref{thm:gm-reduction}\\
辛积分器 & 精确求解修正 $\tilde H=H+O(h^r)$ & \cref{thm:gm-backward-error}\\
\bottomrule
\end{tabular}
\end{table}

\section*{常见误解}
\addcontentsline{toc}{section}{常见误解}
\begin{itemize}
  \item \textbf{「任何合法 $\mathbf C$ 都使 $\dot{\mathbf M}-2\mathbf C$ 反对称」}——$\mathbf C$ 不唯一，只有 Christoffel 构造保证全向量反对称（强版本），Slotine--Li 证明所需（\cref{thm:gm-skew}）。
  \item \textbf{「Euler 方程的 $\boldsymbol\omega\times\mathbf I_c\boldsymbol\omega$ 是物理力矩」}——它是体坐标系旋转的几何效应（$\operatorname{ad}^*$ 项），空间系下无此项（\cref{eq:gm-euler-rigid} 旁）。
  \item \textbf{「Euler--Poincaré 只是 Euler--Lagrange 换记号」}——前者在 $\mathfrak g$（$n$ 维）、后者在 $TG$（$2n$ 维），维度减半（\cref{thm:gm-euler-poincare}）。
  \item \textbf{「math $\operatorname{ad}^*$ 与 Featherstone $\times^*$ 通用」}——二者差一负号，混用得错号方程，本章头号易错点（\cref{eq:gm-ad-se3} 旁）。
  \item \textbf{「保辛 $=$ 保能量」}——辛积分器精确保的是修正 Hamiltonian，能量有界振荡而非守恒（\cref{thm:gm-backward-error}）。
  \item \textbf{「$SE_2(3)=\SE(3)\times\mathbb R^3$」}——作为集合同构但群乘法不同（$\mathbf v_1+\mathbf R_1\mathbf v_2$ vs $\mathbf v_1+\mathbf v_2$），导致 InEKF 误差动力学是否自主的差异（\cref{ch:isam2}）。
\end{itemize}

\section*{延伸阅读}
\addcontentsline{toc}{section}{延伸阅读}
\begin{itemize}
  \item \textbf{几何力学与约化（黄金标准）}：Marsden \& Ratiu《Introduction to Mechanics and Symmetry》——Euler--Poincaré、Lie--Poisson、Marsden--Weinstein 约化的标准参考；Holm, Schmah \& Stoica《Geometric Mechanics and Symmetry》\cite{holm2009geometric}——从 Euler top 到落猫问题的几何相位。
  \item \textbf{辛几何}：da Silva《Lectures on Symplectic Geometry》\cite{dasilva2008lectures}（Darboux/动量映射，纯数学最清晰）、Arnold《Mathematical Methods of Classical Mechanics》\cite{arnold1989mathematical}（Arnold--Liouville/Poincaré 回归/Liouville）。
  \item \textbf{辛积分器}：Hairer, Lubich \& Wanner《Geometric Numerical Integration》\cite{hairer2006geometric}（向后误差分析的标准参考）；冯康、秦孟兆《哈密尔顿系统的辛几何算法》\cite{feng2003symplectic}（国内辛算法标准）。
  \item \textbf{机器人动力学教材}：Spong, Hutchinson \& Vidyasagar《Robot Modeling and Control》\cite{spong2006robot}（Christoffel/$\dot{\mathbf M}-2\mathbf C$/被动性）、Lynch \& Park《Modern Robotics》\cite{lynch2017modern}（$\SE(3)$ 动力学与 (8.40)）、Murray, Li \& Sastry《A Mathematical Introduction to Robotic Manipulation》\cite{murray1994mathematical}、Goldstein 等《Classical Mechanics》\cite{goldstein2002classical}（分析力学物理直觉）。
  \item \textbf{原始文献}：Poincaré 1901（Euler--Poincaré 方程\cite{poincare1901forme}）、Marsden \& Weinstein 1974（辛约化\cite{marsden1974reduction}）、Ruth 1983/Forest--Ruth 1990/Yoshida 1990（辛积分器\cite{ruth1983canonical}\cite{forest1990fourth}\cite{yoshida1990construction}）、Feng Kang 1985（辛差分\cite{feng1985difference}）、Kobilarov--Marsden 2011（DMOC\cite{kobilarov2011discrete}）。
\end{itemize}


